\documentclass[11pt,reqno]{amsart}

\usepackage[a4paper,margin=1in]{geometry}
\usepackage{amsmath,amssymb,amsthm,mathtools}
\usepackage{booktabs}
\usepackage{microtype}
\usepackage{xcolor}
\usepackage[colorlinks=true,linkcolor=blue!60!black,citecolor=green!40!black,urlcolor=blue!60!black]{hyperref}
\usepackage[capitalise,nameinlink]{cleveref}

\numberwithin{equation}{section}

\theoremstyle{plain}
\newtheorem{theorem}{Theorem}[section]
\newtheorem{lemma}[theorem]{Lemma}
\newtheorem{proposition}[theorem]{Proposition}
\newtheorem{corollary}[theorem]{Corollary}
\newtheorem{conjecture}[theorem]{Conjecture}
\theoremstyle{definition}
\newtheorem{definition}[theorem]{Definition}
\theoremstyle{remark}
\newtheorem{remark}[theorem]{Remark}

\crefname{conjecture}{Conjecture}{Conjectures}

\newcommand{\Q}{\mathbb{Q}}
\newcommand{\Z}{\mathbb{Z}}
\newcommand{\C}{\mathbb{C}}
\newcommand{\F}{\mathbb{F}}
\newcommand{\E}{\mathbb{E}}
\newcommand{\one}{\mathbf{1}}
\newcommand{\T}{\mathsf{T}}
\newcommand{\HL}{\mathcal{H}}
\newcommand{\Id}{\mathrm{Id}}
\newcommand{\M}{\mathcal{M}}
\newcommand{\qbin}[2]{\begin{bmatrix}#1\\#2\end{bmatrix}_q}
\newcommand{\diag}{\operatorname{diag}}
\newcommand{\Li}{\operatorname{Li}}
\newcommand{\file}[1]{\nolinkurl{#1}}
\def\UrlBreaks{\do\_\do\/\do\-}

\begin{document}

\title[Modular shifts of tadpole Nahm sums]{Modular shifts of tadpole Nahm sums:\\ classifications, companion identities,\\ and boundary calculus}

% ---- Author information: to be completed by the author(s) ----
\author{[Author name]}
\address{[Department, Institution, Address]}
\email{[email address]}

\subjclass[2020]{Primary 11P84, 11F03; Secondary 33D15, 05E05, 13P10}
\keywords{Nahm sums, tadpole Cartan matrix, modular functions, Rogers--Ramanujan type identities, Hall--Littlewood polynomials, Wronskians, asymptotic expansions, Gr\"obner bases}

\begin{abstract}
For the tadpole matrix $T_r$ we study the set $\M_r$ of rational pairs $(B,C)$ for which the Nahm sum $F_{r,B,C}$ is a modular function. We determine $\M_r$ completely for $r=2,3,4,5,7,9$, where it has $5,6,8,9,12,15$ elements, respectively, and we show that $\M_6$ has at most nine elements; no bound on the size or denominators of $B$ and $C$ is assumed. The modular examples in ranks three, four and five were known, and our contribution there is exhaustiveness. Necessity follows from the vanishing of cumulants in the Vlasenko--Zwegers asymptotic expansion at a cusp, combined with exact elimination; in ranks seven and nine we transfer a finite-field dimension bound to $\Q$ by means of a criterion on highest homogeneous parts. In every odd rank $r=2m+1$ the classifications suggest a family of $3(m+1)$ modular pairs. We prove a companion Wronskian identity for these pairs at the indices $k\in\{0,1,\dots,6,m\}$ using Hall--Littlewood identities of Bartlett and Warnaar together with exact telescoping certificates. Consequently, all proposed pairs are modular through rank fifteen, and every odd rank $r\ge 15$ has at least twenty-four. We also develop a boundary calculus valid for all indices: an explicit finite kernel, an affine folding formula, and a reduction of every boundary shift to the first $m+1$ shifts with coefficients in $\Z[q]$, which form a minimal basis when the Hall--Littlewood variables are free. A polynomial recurrence relates the coefficients to $q$-Fibonacci polynomials. We formulate a local identity that would complete the companion induction, and an even-rank analogue of the odd-rank conjecture, together with even-rank Wronskian identities that would complete the classification in rank six. We prove these identities in one even sector for the first six indices, using a type-B identity of Bartlett and Warnaar, and test the rest up to rank fourteen.
\end{abstract}

\maketitle

\section{Introduction}\label{sec:intro}

\subsection{Nahm sums and the tadpole matrix}
For a positive definite symmetric matrix $A\in\Q^{r\times r}$, a vector $B\in\Q^r$ and a scalar $C\in\Q$, the Nahm sum
\[
f_{A,B,C}(q)=\sum_{n\in\Z_{\ge0}^r}\frac{q^{\frac12 n^\T An+B^\T n+C}}{(q;q)_{n_1}\cdots(q;q)_{n_r}}
\]
is holomorphic for $|q|<1$. Nahm's problem asks for all triples $(A,B,C)$ for which $f_{A,B,C}$ is a modular function \cite{Nahm07,Zagier07}. In rank one, Zagier showed that there are exactly seven such triples \cite{Zagier07}. In higher rank no classification is known. Zagier compiled lists of candidates in ranks two and three, and Vlasenko and Zwegers developed the asymptotic analysis at cusps that underlies much of the subsequent work, including counterexamples to the most direct form of Nahm's conjecture \cite{VZ11}; one direction of the conjecture was later proved by Calegari, Garoufalidis and Zagier \cite{CGZ}. Asymptotic conditions of this kind have also been used to search for modular companions of known identities \cite{Kanade19}.

This paper concerns the \emph{tadpole} family. For $r\ge1$, let $T_r$ be the symmetric matrix with diagonal entries $2,\dots,2,1$, entries $-1$ immediately above and below the diagonal, and all other entries zero. For $B\in\Q^r$, $C\in\Q$ and $q=e^{2\pi i\tau}$ with $\operatorname{Im}\tau>0$, put
\begin{equation}\label{eq:nahm}
F_{r,B,C}(q)=\sum_{n\in\Z_{\ge0}^r}\frac{q^{\frac12 n^\T T_rn+B^\T n+C}}{\prod_{j=1}^r(q;q)_{n_j}},\qquad (q;q)_m=\prod_{j=1}^m(1-q^j),
\end{equation}
where an empty product is $1$ and rational powers mean $q^a=e^{2\pi ia\tau}$. Since $T_r$ is positive definite, the series converges on the upper half-plane. Calinescu, Milas and Penn conjectured that the principal tadpole sums, which arise from principal subspaces of the twisted affine algebra $A_{2r}^{(2)}$, are modular \cite{CMP16}. The conjecture was proved by Milas and Wang for $r=3$, together with six modular shifts \cite{MW24}, and by Shi and Wang for $r=4,5$, together with principal and companion sums \cite{SW26}. Further rank-four companions were obtained in \cite{SW26b}; three of these identities were proved independently by Li \cite{Li26}, and related rank-four families arise from lift-dual operations \cite{CW25,CW26}. Chern, Tran and Wakhare showed that the tadpole Nahm sum is a specialization of a series of Bartlett and Warnaar \cite{BW15}, and obtained a Wronskian formula for the principal and last-coordinate half-shifted sums in every rank \cite{CTW26}.

We study the corresponding classification problem: determine
\[
\M_r=\{(B,C)\in\Q^r\times\Q:\ F_{r,B,C}\ \text{is a modular function}\}.
\]
Here a modular function is a function invariant under a congruence subgroup of $\mathrm{SL}_2(\Z)$ and meromorphic at its cusps. Proving that particular pairs lie in $\M_r$ (sufficiency) and proving that a given list exhausts $\M_r$ (necessity) are different problems. Most of the literature cited above concerns sufficiency.

\subsection{Classification results}
The odd-rank classifications obtained below fit a single family.

\begin{conjecture}[Odd-rank classification]\label{conj:odd}
For $r=2m+1$ with $m\ge1$, the set $\M_r$ consists of the following $3(m+1)$ pairs. Choose
\[
z\in\{-\tfrac12,0,\tfrac12\},\qquad 0\le x_1\le\cdots\le x_m\le1,\qquad x_j\in\Z,
\]
put $k=x_1+\cdots+x_m$, and set
\begin{align*}
B&=(2z-x_1,\,x_1-2z,\,\dots,\,2z-x_m,\,x_m-2z,\,z)^\T,\\
C&=-\frac{r(2r+1)}{48(r+2)}+\frac{k(r+1-k)}{2(r+2)}-kz+\frac{2r-1}{4}z^2.
\end{align*}
\end{conjecture}

The scalar satisfies $C(k,-\tfrac12)-C(k,\tfrac12)=k$, as required by the identity \eqref{eq:halfshift} below. For $m=0$ the formula gives $B=(z)$ and $C=-\tfrac1{48}+\tfrac{z^2}4$, which are exactly the three rank-one modular triples with $A=1$ in Zagier's classification \cite{Zagier07}.

\begin{theorem}\label{thm:main-class}
For $r\in\{2,3,4,5,7,9\}$ the set $\M_r$ is finite and has the following cardinality.
\begin{center}
\begin{tabular}{ccc}
\toprule
Rank $r$ & $|\M_r|$ & Statement\\
\midrule
2 & 5 & \cref{thm:r2}\\
3 & 6 & \cref{thm:r3}\\
4 & 8 & \cref{thm:r4}\\
5 & 9 & \cref{thm:r5}\\
7 & 12 & \cref{cor:r7}\\
9 & 15 & \cref{thm:r9}\\
\bottomrule
\end{tabular}
\end{center}
In particular, \cref{conj:odd} holds for $r=3,5,7,9$. Moreover, $\M_6$ has at most nine elements (\cref{thm:r6}). No bound is imposed on the denominators or sizes of the rational shifts.
\end{theorem}

For $r=3,4,5$ every element of $\M_r$ was already known to be modular \cite{MW24,SW26,SW26b}; the new statement is that the lists are complete. For $r=2,7,9$ both directions are proved here; in rank two the five modular sums are infinite products given by the Göllnitz--Gordon and Lebesgue identities. The even ranks two, four and six suggest a second family, with $2m+3$ pairs in rank $2m$ for $m\ne2$ (\cref{conj:even}). It agrees with the cusp conditions in ranks eight and ten, and \cref{conj:odd} agrees with them in ranks eleven and thirteen (\cref{rem:highrank}). The even-rank sums also appear to satisfy Wronskian identities with an alternating theta character (\cref{conj:evenwron}); we have checked them as exact $q$-series in all even ranks up to fourteen, and they would imply the sufficiency direction of \cref{conj:even}, in particular completing the classification in rank six (\cref{prop:evenwron}). Using the type-B identity of Bartlett and Warnaar, we prove the first companions in one even sector for $1\le k\le\min(6,m-1)$ (\cref{thm:evencomp}); in rank six this proves four of the nine pairs. The proofs of necessity are computer-assisted; see \cref{subsec:computation}.

\subsection{Companion identities}
For $r=2m+1$ and $0\le k\le m$, let $B(k,z)$ and $C(k,z)$ denote the pair of \cref{conj:odd} with $x_1=\cdots=x_{m-k}=0$ and $x_{m-k+1}=\cdots=x_m=1$. In \cref{sec:theta} we propose an explicit Wronskian (theta-determinant) formula for $F_{r,B(k,z),0}$ (\cref{conj:wron}), and show that it implies modularity with exactly the scalar $C(k,z)$ (\cref{prop:rhs-modular}).

\begin{theorem}\label{thm:main-comp}
Let $r=2m+1$. The identity \eqref{eq:wron} holds for every
\[
k\in\{0,1,2,3,4,5,6,m\}\cap\{0,\dots,m\},\qquad z\in\{0,\tfrac12\}.
\]
Consequently $(B(k,z),C(k,z))\in\M_r$ for these $k$ and all $z\in\{-\frac12,0,\frac12\}$. In particular, every pair of \cref{conj:odd} is modular for $r\le15$, and $\M_r$ contains at least twenty-four elements for every odd $r\ge15$.
\end{theorem}

The cases $k\in\{0,1,m\}$ (\cref{thm:basic}) follow from two identities of Bartlett and Warnaar \cite{BW15} through a finite determinant calculation. The case $k=2$ (\cref{thm:k2}) uses a short telescoping certificate, which is printed in full, and the cases $3\le k\le6$ (\cref{thm:k3to6}) use larger certificates of the same type. Two of the identities, $(k,z)=(0,0)$ and $(m,\tfrac12)$, are the formulas of Chern, Tran and Wakhare \cite[Theorem~1.4]{CTW26}.

\subsection{Boundary calculus}
\Cref{sec:boundary} treats the Hall--Littlewood side of the companion problem for arbitrary indices. We give an explicit finite kernel for the relevant products of tridiagonal matrices (\cref{prop:kernel}), an affine folding formula for the boundary inverse (\cref{thm:folding}), and a reduction of every boundary shift to the first $m+1$ shifts with coefficients in $\Z[q]$ (\cref{thm:basis,cor:reduction}). These $m+1$ functions form a minimal basis when the Hall--Littlewood variables are free. A polynomial recurrence (\cref{prop:recurrence}) connects the coefficients with $q$-Fibonacci polynomials \cite{Cigler03}. The only missing ingredient for a proof of \cref{conj:wron} at all indices is a reduction on the Nahm side; \cref{sec:open} isolates it as an explicit local identity (\cref{prop:induction}).

\subsection{Relation to previous work}
The six rank-three identities and their modularity are due to Milas and Wang \cite{MW24}; the nine rank-five pairs and the first four rank-four pairs are due to Shi and Wang \cite{SW26}, and the remaining four rank-four pairs to \cite{SW26b} (see also \cite{Li26}). Up to a permutation of coordinates, the rank-four candidate set already appears in \cite[Table~3]{CW25}. The simultaneous-shift relation (\cref{lem:simshift}) is a consequence of the rank reduction of Shi and Wang, and the coefficient inversion of \cref{lem:inverse} is a rescaled Bailey inversion. The theta numerators in \eqref{eq:wron} are of standard affine-character form, so their modularity is not in itself a new statement about Nahm sums. To our knowledge, the exhaustive classifications of \cref{thm:main-class}, the companion identities of \cref{thm:main-comp} other than the two cases $(k,z)=(0,0)$ and $(m,\frac12)$ of \cite{CTW26}, and the boundary calculus of \cref{sec:boundary} are new.

\subsection{Computations}\label{subsec:computation}
The proofs of necessity and of the higher companion identities use three kinds of computation.
\begin{enumerate}
\item In ranks three, four and five, the eliminations are proved by rational polynomial certificates (explicit ideal-membership cofactors and Gr\"obner bases), which are checked in exact rational arithmetic.
\item In ranks two, six, seven and nine, an affine elimination over a prime field and a computation with highest homogeneous parts are combined through \cref{lem:transfer}. The cusp polynomials are computed in exact number-field arithmetic.
\item The companion identities for $3\le k\le6$ are proved by Laurent-polynomial telescoping identities, which are checked by exact integer expansion.
\end{enumerate}
\Cref{app:verification} describes these computations. Numerical $q$-series comparisons appear only in \cref{sec:open}, as evidence for conjectures, and are not used in any proof.

\subsection{Organization}
\Cref{sec:cusp} sets up the cusp obstructions and the transfer lemma. \Cref{sec:low} proves the classifications. \Cref{sec:theta} introduces the Wronskian identities and proves the basic companions, \cref{sec:boundary} develops the boundary calculus, and \cref{sec:cert} proves the companions with $2\le k\le6$. \Cref{sec:open} states the remaining problems. \Cref{app:data} records the obstruction data in ranks three, four and five.

\subsection{Notation}
We write $(a;q)_n=\prod_{s=0}^{n-1}(1-aq^s)$ for $n\ge0$, $(a;q)_\infty=\prod_{s\ge0}(1-aq^s)$, and $(a_1,\dots,a_t;q)_n=\prod_{i=1}^t(a_i;q)_n$, also for $n=\infty$. For a nonnegative integer $a$, the Gaussian binomial coefficient is
\[
\qbin{a}{b}=\frac{(q;q)_a}{(q;q)_b(q;q)_{a-b}}\qquad(0\le b\le a),
\]
and is zero otherwise. We write $e_i$ for the $i$th coordinate vector, $\one=(1,\dots,1)$, $\delta_{ij}$ for the Kronecker delta, and $\eta(\tau)=q^{1/24}\prod_{n\ge1}(1-q^n)$ for the Dedekind eta function.

\section{Cusp obstructions and specialization}\label{sec:cusp}

\subsection{Cusp coefficients}
We use the asymptotic theorem of Vlasenko and Zwegers \cite[Theorem~2.3 and Corollary~3.2]{VZ11}; theorem numbering refers to the arXiv version. Write $F_{A,B,C}$ for \eqref{eq:nahm} with $T_r$ replaced by a positive definite rational matrix $A$. Let $Q_j\in(0,1)$ solve
\[
1-Q_i=\prod_{j=1}^rQ_j^{A_{ij}},\qquad \xi_i=\frac{Q_i}{1-Q_i},\qquad M=(A+\diag\xi)^{-1}.
\]
There is a unique such solution. With $q=e^{-\varepsilon}$ and $\varepsilon\downarrow0$, their expansion has the form
\begin{equation}\label{eq:asymp}
F_{A,B,C}(e^{-\varepsilon})\sim e^{\alpha/\varepsilon}\beta e^{-(C+\gamma_0)\varepsilon}\bigl(1+c_1\varepsilon+c_2\varepsilon^2+c_3\varepsilon^3+\cdots\bigr),
\end{equation}
where $\beta>0$, $\gamma_0=r/24+\sum_i\xi_i/12$, and $\alpha=\sum_i(L(1)-L(Q_i))$. Here $L(x)=\Li_2(x)+\frac12\log x\log(1-x)$ is the Rogers dilogarithm and $\Li_2(x)=\sum_{n\ge1}x^n/n^2$ for $0<x<1$. The constants $c_j$ are Gaussian moments specified below.

At a cusp, a nonzero modular function has an exponential leading term and exponentially small relative corrections. Comparing this with \eqref{eq:asymp} gives the necessary conditions
\begin{equation}\label{eq:cuspcond}
C=c_1-\gamma_0,\qquad \kappa_2:=c_2-\tfrac12c_1^2=0,\qquad \kappa_3:=c_3-c_1c_2+\tfrac13c_1^3=0.
\end{equation}
The same comparison shows that a modular Nahm sum in this setting has weight zero. Conditions \eqref{eq:cuspcond} are only necessary; we make no sufficiency claim for truncated cusp expansions.

We now give the finite procedure used to compute the coefficients, for general rank $r$ and truncation order $N$. Let $X=(X_1,\dots,X_r)$ be a centred real Gaussian vector with covariance matrix $M$. Define polynomials
\[
P_1(x)=x,\qquad P_{n+1}(x)=x(1+x)P_n'(x),
\]
so that $P_n(\xi_i)=\Li_{1-n}(Q_i)$. Let $\mathcal B_j$ be the Bernoulli numbers defined by $t/(e^t-1)=\sum_{j\ge0}\mathcal B_jt^j/j!$, so $\mathcal B_1=-\frac12$. For $1\le k\le 2N$, put
\[
H_k(X)=\one_{k=1}\sum_{i=1}^r\Bigl(b_i+\frac{\xi_i}2\Bigr)X_i-\sum_{i=1}^r\sum_{\substack{g\ge3,\ 0\le j\le g\\ g+j-2=k}}\frac{\mathcal B_j}{j!\,(g-j)!}P_{g-1}(\xi_i)X_i^{g-j},
\]
where $\one_{k=1}$ is $1$ for $k=1$ and $0$ otherwise. Define $D_0=1$ and
\[
D_\ell=\frac1\ell\sum_{k=1}^\ell kH_kD_{\ell-k}\quad(1\le \ell\le2N),\qquad c_j=\E[D_{2j}(X)]\quad(1\le j\le N).
\]
This is coefficient extraction from $\exp\bigl(\sum_{k=1}^{2N}H_kh^k\bigr)$ with $h=\sqrt\varepsilon$.

The Gaussian moments are determined without integration. If $\mu(a)=\E\bigl[\prod_{j=1}^rX_j^{a_j}\bigr]$, then $\mu(0)=1$, moments of odd total degree vanish, and, whenever $a_i>0$,
\[
\mu(a)=\sum_{j=1}^r(a_j-\delta_{ij})M_{ij}\,\mu(a-e_i-e_j),
\]
where terms with zero prefactor are omitted.

\begin{remark}\label{rem:cusp-check}
The procedure above is the expansion of \cite[Theorem~2.3]{VZ11} written in moment form. In rank three, the resulting saddle point and first-order condition agree with those computed in \cite[Section~4.1]{MW24} (see \cref{app:r3}).
\end{remark}

\subsection{Degree control}

\begin{lemma}[Degree of a logarithmic cusp coefficient]\label{lem:degree}
In every rank, the coefficient $\kappa_j(B)$ of $h^{2j}$ in
\[
\log\E\Bigl[\exp\Bigl(\sum_{a\ge1}h^aH_a(X)\Bigr)\Bigr]
\]
has total degree at most $j+1$ in $B$. In particular, $c_1-\gamma_0$, $\kappa_2$ and $\kappa_3$ have degrees at most two, three and four, respectively.
\end{lemma}

\begin{proof}
The only $B$-dependent term in the exponent is $h\sum_ib_iX_i$. By the definition of $H_a$, every other monomial of order $h^a$ has Gaussian degree at most $a+2$. Expand a Gaussian moment as a sum over pairings, drawing one edge for each pair of Gaussian factors. Taking the logarithm cancels the disconnected pairing graphs, by the moment--cumulant formula.

Consider a surviving connected graph with $\ell$ vertices from $h\sum_ib_iX_i$ and $s$ other vertices, of orders $a_1,\dots,a_s$ and Gaussian degrees $d_1,\dots,d_s$. It has $(\ell+\sum_id_i)/2$ edges, and since it is connected this number is at least $\ell+s-1$. Since $d_i\le a_i+2$,
\[
\ell\le\sum_i(d_i-2)+2\le\sum_ia_i+2.
\]
The graph contributes to $h^{2j}$ only if $\ell+\sum_ia_i=2j$, and therefore $\ell\le j+1$. The degree in $B$ is $\ell$. Subtracting the $B$-independent quantity $\gamma_0$ does not affect the first-order bound.
\end{proof}

\subsection{Transfer from a finite field}
The following criterion replaces an elimination over $\Q$ by one over a prime field. Its hypothesis on highest homogeneous parts excludes the loss of solutions at infinity; an affine calculation at a single prime would not suffice.

\begin{lemma}[Finite specialization]\label{lem:transfer}
Let $f_1,\dots,f_s\in\Z[X_1,\dots,X_n]$ be nonzero polynomials, and let $f_i^{\mathrm{top}}$ be the homogeneous part of $f_i$ of highest total degree. Let $p$ be a prime, and let a bar denote reduction modulo $p$. Write $X=(X_1,\dots,X_n)$ and suppose that
\[
\F_p[X]/\bigl(\overline{f_1^{\mathrm{top}}},\dots,\overline{f_s^{\mathrm{top}}}\bigr)\quad\text{is finite-dimensional.}
\]
If $d=\dim_{\F_p}\F_p[X]/(\bar f_1,\dots,\bar f_s)$, then
\[
\dim_\Q\Q[X]/(f_1,\dots,f_s)\le d.
\]
\end{lemma}

\begin{proof}
Put $R=\Z_{(p)}$, the ring of rational numbers whose denominators are not divisible by $p$. By hypothesis, for some $N$ the degree-$N$ monomial multiples of the $\overline{f_i^{\mathrm{top}}}$ span all homogeneous polynomials of degree $N$ over $\F_p$. The corresponding integer coefficient matrix therefore has a maximal minor that is a unit in $R$, so every monomial of degree $N$ is an $R$-linear combination of monomial multiples of the $f_i^{\mathrm{top}}$. Replacing each top part by $f_i$ expresses that monomial, modulo $(f_1,\dots,f_s)$, in terms of monomials of smaller degree. By induction,
\[
A=R[X]/(f_1,\dots,f_s)
\]
is generated as an $R$-module by the monomials of degree less than $N$. Choose lifts to $A$ of a basis of $A/pA\cong\F_p[X]/(\bar f_1,\dots,\bar f_s)$. By Nakayama's lemma these $d$ elements generate $A$ over the local ring $R$. Tensoring with $\Q$ proves the bound.
\end{proof}

\section{Classifications in low rank}\label{sec:low}

\subsection{Rank two}

\begin{theorem}[Rank-two classification]\label{thm:r2}
Let $B=(b_1,b_2)^\T\in\Q^2$ and $C\in\Q$. Then $F_{2,B,C}$ is a modular function if and only if $(B,C)$ is one of the following five pairs, for which, with $w=q^{1/2}$,
\begin{center}
\begin{tabular}{rrrl}
\toprule
$b_1$ & $b_2$ & $C$ & $F_{2,B,0}(q)$\\
\midrule
$0$ & $0$ & $-5/96$ & $(-w;q)_\infty\big/\bigl((w;w^8)_\infty(w^4;w^8)_\infty(w^7;w^8)_\infty\bigr)$\\
$1$ & $0$ & $19/96$ & $(-w;q)_\infty\big/\bigl((w^3;w^8)_\infty(w^4;w^8)_\infty(w^5;w^8)_\infty\bigr)$\\
$0$ & $1/2$ & $1/24$ & $(-q;q)_\infty^2(-q;q^2)_\infty$\\
$1$ & $-1/2$ & $1/6$ & $2(-q;q)_\infty^2(-q^2;q^2)_\infty$\\
$2$ & $-3/2$ & $25/24$ & $2q^{-1}(-q;q)_\infty^2(-q;q^2)_\infty$\\
\bottomrule
\end{tabular}
\end{center}
\end{theorem}

\begin{proof}
\emph{Sufficiency.} Summing over $n_2$ with Euler's identity $\sum_{n\ge0}q^{n(n-1)/2}z^n/(q;q)_n=(-z;q)_\infty$ gives
\begin{equation}\label{eq:r2euler}
F_{2,B,0}(q)=\sum_{n\ge0}\frac{q^{n^2+b_1n}}{(q;q)_n}\,(-q^{1/2+b_2-n};q)_\infty .
\end{equation}
If $b_2=0$, then $(-q^{1/2-n};q)_\infty=q^{-n^2/2}(-w;q)_n(-w;q)_\infty$, so
\[
F_{2,B,0}=(-w;q)_\infty\sum_{n\ge0}\frac{w^{n^2+2b_1n}(-w;w^2)_n}{(w^2;w^2)_n}.
\] For $b_1=0$ and $b_1=1$ these are the two Göllnitz--Gordon identities \cite{Gordon65,Gollnitz67} in the variable $w$, which give the first two rows. If $b_2=\frac12-s$ with $s\in\{0,1,2\}$, then $(-q^{1-s-n};q)_\infty=q^{-\binom{n+s}2}(-1;q)_{n+s}(-q;q)_\infty$, and in the three cases $(b_1,s)=(0,0),(1,1),(2,2)$ the sum \eqref{eq:r2euler} becomes
\[
\sum_{n\ge0}\frac{q^{\binom{n+1}2}(-1;q)_n}{(q;q)_n},\qquad 2\sum_{n\ge0}\frac{q^{\binom{n+1}2}(-q;q)_n}{(q;q)_n},\qquad 2(1+q)q^{-1}\sum_{n\ge0}\frac{q^{\binom{n+1}2}(-q^2;q)_n}{(q;q)_n},
\]
each multiplied by $(-q;q)_\infty$. Lebesgue's identity $\sum_{n\ge0}(a;q)_nq^{\binom{n+1}2}/(q;q)_n=(aq;q^2)_\infty(-q;q)_\infty$ \cite{Andrews76} with $a=-1,-q,-q^2$ gives the last three rows. Every product in the table is a quotient of (generalized) Dedekind eta functions in $\tau$ or $\tau/2$, so it becomes a modular function on a congruence subgroup after multiplication by the power of $q$ that makes its order at $i\infty$ match its eta-quotient normalization; for the five rows this power is exactly the displayed $C$.

\emph{Necessity.} Here $t=1-\cos(\pi/4)=1-\frac{\sqrt2}2$, and the saddle point and covariance lie in $\Q(t)=\Q(\sqrt2)$. As in \cref{thm:r9}, take the irrational component of $c_1-\gamma_0$ and both components of $\kappa_2$ and $\kappa_3$; this gives five primitive integer polynomials in $(b_1,b_2)$, of degrees $2,3,3,4,4$. Modulo $p=1000003$ their ideal has a quotient of dimension five, and the ideal of their highest homogeneous parts has a finite quotient. \Cref{lem:transfer} bounds the rational quotient dimension by five. The five pairs in the table are distinct rational zeros, so they are all the zeros, and the first-order condition gives the displayed $C$.
\end{proof}

The last row is twice $q^{-1}$ times the third. This is the rank-two instance of an even-rank half-shift relation; see \cref{conj:even}.

\subsection{Rank three}

\begin{theorem}[Rank-three classification]\label{thm:r3}
Let $B=(b_1,b_2,b_3)^\T\in\Q^3$ and $C\in\Q$. Then $F_{3,B,C}$ is a modular function if and only if $(B,C)$ is one of the following pairs.
\begin{center}
\begin{tabular}{rrrr}
\toprule
$b_1$ & $b_2$ & $b_3$ & $C$\\
\midrule
$0$ & $0$ & $0$ & $-7/80$\\
$0$ & $0$ & $1/2$ & $1/40$\\
$1$ & $-1$ & $1/2$ & $9/40$\\
$-1$ & $1$ & $0$ & $17/80$\\
$-1$ & $1$ & $-1/2$ & $9/40$\\
$-2$ & $2$ & $-1/2$ & $41/40$\\
\bottomrule
\end{tabular}
\end{center}
\end{theorem}

The sufficiency of these six pairs is due to Milas and Wang \cite[Theorems~1.2--1.3]{MW24}. \Cref{app:r3} supplies the exact necessity certificate.

The classification is especially simple in the coordinates
\begin{equation}\label{eq:r3coords}
u=b_1+b_2,\qquad v=b_2+2b_3,\qquad z=b_3.
\end{equation}
The allowed shifts are precisely those with
\begin{equation}\label{eq:r3zeros}
u=0,\qquad v(v-1)=0,\qquad z(4z^2-1)=0.
\end{equation}

\begin{proof}[Proof of \cref{thm:r3}]
The asymptotic conditions imply $E_0=\cdots=E_4=0$ for the polynomials defined in \cref{app:r3}. By \cref{lem:r3elim}, $u=0$, $v\in\{0,1\}$ and $z\in\{-\frac12,0,\frac12\}$. The inverse change of coordinates is
\[
b_1=u-v+2z,\qquad b_2=v-2z,\qquad b_3=z,
\]
and these six choices give the six rows of the table. Substituting them into $C_0(B)$ from \eqref{eq:r3first} gives the displayed scalars. Their modularity is established in \cite{MW24}.
\end{proof}

\subsection{Rank four}

\begin{theorem}[Rank-four classification]\label{thm:r4}
Let $B=(b_1,b_2,b_3,b_4)^\T\in\Q^4$ and $C\in\Q$. Then $F_{4,B,C}$ is a modular function if and only if $(B,C)$ is one of the following eight pairs.
\begin{center}
\begin{tabular}{rrrrrl}
\toprule
$b_1$ & $b_2$ & $b_3$ & $b_4$ & $C$ & Sufficiency\\
\midrule
$0$ & $0$ & $0$ & $0$ & $-1/8$ & \cite[(1.20)]{SW26}\\
$0$ & $-1$ & $1$ & $0$ & $5/24$ & \cite[(1.21)]{SW26}\\
$0$ & $0$ & $0$ & $1/2$ & $0$ & \cite[(1.22)]{SW26}\\
$1$ & $-1$ & $1$ & $0$ & $3/8$ & \cite[(1.23)]{SW26}\\
$1$ & $-1$ & $1$ & $-1/2$ & $1/3$ & \cite[(1.10), Theorem~1.6]{SW26b}\\
$2$ & $-2$ & $1$ & $-1/2$ & $5/4$ & \cite[(1.12), Theorem~1.6]{SW26b}\\
$3/2$ & $-3/2$ & $1/2$ & $0$ & $9/16$ & \cite[(1.13), Theorem~1.6]{SW26b}\\
$1$ & $-1$ & $0$ & $1/2$ & $1/4$ & \cite[(1.14), Theorem~1.6]{SW26b}\\
\bottomrule
\end{tabular}
\end{center}
\end{theorem}

All eight modular examples are known; after a permutation of coordinates the candidate set appears in \cite[Table~3 and the subsequent discussion of $B_8'$]{CW25}. The identities (1.12)--(1.14) of \cite{SW26b} were also proved independently in \cite{Li26}. The content of \cref{thm:r4} is that the list is exhaustive among rational $B$ and $C$.

\begin{proof}[Proof of \cref{thm:r4}]
Modularity and rationality give the seven obstructions of \cref{app:r4}. By \cref{lem:r4elim}, $B$ is one of the eight listed vectors, and the first-order condition gives the displayed values $C=C_4(B)$. Sufficiency follows from the cited identities. For the last four rows the normalized modular functions are, respectively,
\[
6\,\frac{\eta(2\tau)^2\eta(3\tau)^3}{\eta(\tau)^5},\qquad 8\,\frac{\eta(2\tau)^6}{\eta(\tau)^6},\qquad 2\,\frac{\eta(\tau)^3}{\eta(\tau/2)^3},\qquad 4\,\frac{\eta(2\tau)^6}{\eta(\tau)^6};
\]
these are consequences of \cite[(1.10), (1.12)--(1.14), Theorem~1.6]{SW26b}.
\end{proof}

In rank four the fourth cusp coefficient is needed to obtain a radical ideal. We do not claim that this number of coefficients is minimal.

\subsection{Rank five}

\begin{theorem}[Rank-five classification]\label{thm:r5}
Let $B=(b_1,\dots,b_5)^\T\in\Q^5$ and $C\in\Q$. Define
\[
u=b_1+b_2,\quad v=b_3+b_4,\quad x=b_2+2b_5,\quad y=b_4+2b_5,\quad z=b_5.
\]
Then $F_{5,B,C}$ is a modular function if and only if
\begin{equation}\label{eq:r5set}
u=v=0,\qquad (x,y)\in\{(0,0),(0,1),(1,1)\},\qquad z\in\{-\tfrac12,0,\tfrac12\},
\end{equation}
and
\begin{equation}\label{eq:r5C}
C=\frac{72x+120y-336z(x+y)+756z^2-55}{336}.
\end{equation}
There are exactly nine such pairs $(B,C)$.
\end{theorem}

The inverse coordinates are $B=(u-x+2z,\ x-2z,\ v-y+2z,\ y-2z,\ z)^\T$. For comparison with \cite[Theorem~5.2]{SW26}, write $B=(a,-a,b,-b,c)^\T$. The nine pairs are as follows.
\begin{center}
\begin{tabular}{rrrrl}
\toprule
$a$ & $b$ & $c$ & $C$ & Sufficiency\\
\midrule
$0$ & $0$ & $0$ & $-55/336$ & \cite[(5.29)]{SW26}\\
$1$ & $1$ & $1/2$ & $67/168$ & \cite[(5.30)]{SW26}\\
$-1$ & $-1$ & $0$ & $137/336$ & \cite[(5.31)]{SW26}\\
$0$ & $0$ & $1/2$ & $-5/168$ & \cite[(5.32)]{SW26}\\
$0$ & $-1$ & $0$ & $65/336$ & \cite[(5.33)]{SW26}\\
$1$ & $0$ & $1/2$ & $43/168$ & \cite[(5.34)]{SW26}\\
$-1$ & $-1$ & $-1/2$ & $67/168$ & \cite[(5.35)]{SW26}\\
$-2$ & $-2$ & $-1/2$ & $331/168$ & \cite[(5.36)]{SW26}\\
$-1$ & $-2$ & $-1/2$ & $211/168$ & \cite[(5.37)]{SW26}\\
\bottomrule
\end{tabular}
\end{center}
Shi and Wang use $S(a,b,c;q)=F_{5,(a,-a,b,-b,c),0}(q^4)$, so their scalar exponents must be divided by four in our normalization. All nine modularity statements are proved in \cite{SW26}; the new content of \cref{thm:r5} is exhaustiveness.

\begin{proof}[Proof of \cref{thm:r5}]
The cusp conditions and rationality give the eight polynomial obstructions of \cref{app:r5}, and \cref{lem:r5elim} gives \eqref{eq:r5set}. Reducing the first-order scalar polynomial $C_0$ of \cref{app:r5} modulo this ideal gives \eqref{eq:r5C}. The nine resulting pairs are those in the table, whose modularity is proved in \cite[Theorem~5.2]{SW26}.
\end{proof}

\subsection{Rank six}

\begin{theorem}[Rank-six necessity]\label{thm:r6}
Let $B\in\Q^6$ and $C\in\Q$. If $F_{6,B,C}$ is a modular function, then $(B,C)$ is one of the following nine pairs.
\begin{center}
\begin{tabular}{rrrrrrr}
\toprule
$b_1$ & $b_2$ & $b_3$ & $b_4$ & $b_5$ & $b_6$ & $C$\\
\midrule
$0$ & $0$ & $0$ & $0$ & $0$ & $0$ & $-13/64$\\
$0$ & $0$ & $0$ & $-1$ & $1$ & $0$ & $11/64$\\
$0$ & $-1$ & $1$ & $-1$ & $1$ & $0$ & $27/64$\\
$1$ & $-1$ & $1$ & $-1$ & $1$ & $0$ & $35/64$\\
$0$ & $0$ & $0$ & $0$ & $0$ & $1/2$ & $-1/16$\\
$1$ & $-1$ & $0$ & $0$ & $0$ & $1/2$ & $1/4$\\
$1$ & $-1$ & $1$ & $-1$ & $0$ & $1/2$ & $7/16$\\
$1$ & $-1$ & $1$ & $-1$ & $1$ & $-1/2$ & $1/2$\\
$2$ & $-2$ & $1$ & $-1$ & $1$ & $-1/2$ & $23/16$\\
\bottomrule
\end{tabular}
\end{center}
The first and fifth pairs are modular by \cite{CTW26}, and the second and third by \cref{thm:evencomp}.
\end{theorem}

\begin{proof}
Here $t=1-\cos(\pi/8)$ has degree four over $\Q$. Taking the three irrational components of $c_1-\gamma_0$ and all four components of $\kappa_2$ and $\kappa_3$ gives eleven primitive integer polynomials in $B$, of degrees $2,2,2,3,3,3,3,4,4,4,4$. Modulo $p=1000003$ their ideal has a quotient of dimension nine, and the ideal of their highest homogeneous parts has a finite quotient (of dimension $103$). By \cref{lem:transfer} the rational quotient has dimension at most nine. The nine listed vectors are distinct rational zeros, so the ideal is radical with exactly these zeros, and the first-order condition gives the displayed $C$. The modularity of the first and fifth pairs is the principal and last-coordinate half-shift case of \cite{CTW26}; that of the second and third is \cref{cor:r6more}.
\end{proof}

None of the series other than the first and fifth is an infinite product; their leading coefficients are $5,9,5,6,15,20,30$. As $q$-series to order $q^{24}$, the last series equals $2q^{-1}$ times the seventh, consistently with $C$. We expect all nine pairs to be modular: by \cref{prop:evenwron}, the remaining five would follow from the even-rank Wronskian identities of \cref{conj:evenwron}.

\subsection{Rank seven}

\begin{theorem}[Rank-seven necessity]\label{thm:r7nec}
Let $B\in\Q^7$ and $C\in\Q$. If $F_{7,B,C}$ is a modular function, then $(B,C)$ is one of the twelve pairs of \cref{conj:odd} with $m=3$.
\end{theorem}

\begin{proof}
Let $t=1-\cos(\pi/9)$, the root in $(1/20,7/100)$ of
\[
8t^3-24t^2+18t-1=0.
\]
The saddle coordinates and the covariance matrix lie in $\Q(t)$. Compute the first three cusp coefficients, and take the two irrational components of $c_1-\gamma_0$ and all three components of each of $\kappa_2$ and $\kappa_3$. This gives eight rational obstruction polynomials. Clear denominators and remove the integer content of each; the resulting integer polynomials are recorded in the accompanying data (\cref{app:verification}). Make the invertible change of coordinates
\[
z=b_7,\qquad u_j=b_{2j-1}+b_{2j},\qquad x_j=b_{2j}+2z\qquad(1\le j\le3),
\]
so that $b_{2j-1}=u_j+2z-x_j$ and $b_{2j}=x_j-2z$. Let $I$ be the ideal generated by the eight transformed polynomials in $\Q[u_1,u_2,u_3,x_1,x_2,x_3,z]$. We prove that $I$ equals
\begin{equation}\label{eq:r7ideal}
J=\bigl(u_i,\ x_i^2-x_i\ (1\le i\le3),\ x_ix_j-x_i\ (1\le i<j\le3),\ 4z^3-z\bigr).
\end{equation}

Take $p=1000003$, and let $\bar I$ and $\bar J$ be generated by the reductions of the chosen integer generators of $I$ and $J$. An exact finite-field calculation in \textsc{Singular} \cite{Singular} gives $\bar I=\bar J$, with quotient of dimension twelve. For the top homogeneous parts of the same eight polynomials, the quotient has Hilbert series
\[
1+7s+26s^2+67s^3+131s^4+202s^5+247s^6+231s^7+141s^8,
\]
where $s$ is a formal grading variable. In particular, every monomial of degree nine lies in the ideal generated by the top parts. \Cref{lem:transfer} therefore gives $\dim_\Q\Q[u,x,z]/I\le12$, where $u=(u_1,u_2,u_3)$ and $x=(x_1,x_2,x_3)$.

The twelve rational zeros of $J$ satisfy all eight obstructions, by exact substitution. Their evaluation map is surjective onto $\Q^{12}$; equivalently, the evaluation matrix of the monomials $\{1,x_1,x_2,x_3\}\cdot\{1,z,z^2\}$ is invertible. Hence $\Q[u,x,z]/I$ has dimension at least twelve. Equality follows, the evaluation map is an isomorphism, and $I=J$ is radical. Its twelve zeros have $u_1=u_2=u_3=0$, $0\le x_1\le x_2\le x_3\le1$, $x_i\in\Z$ and $z\in\{-\frac12,0,\frac12\}$. Writing $k=x_1+x_2+x_3$, the first-order condition reduces to
\[
C=-\frac{35}{144}+\frac{k(8-k)}{18}-kz+\frac{13}4z^2,
\]
which is the specialization of \cref{conj:odd}.
\end{proof}

\begin{corollary}[Rank-seven classification]\label{cor:r7}
Let $B\in\Q^7$ and $C\in\Q$. Then $F_{7,B,C}$ is a modular function if and only if $(B,C)$ is one of the twelve pairs of \cref{conj:odd} with $m=3$.
\end{corollary}

\begin{proof}
Necessity is \cref{thm:r7nec}. \Cref{thm:basic,thm:k2} give sufficiency for $k=0,1,2,3$ in all three sectors.
\end{proof}

\subsection{Rank nine}
In this rank it is efficient to derive the cusp conditions directly over a split prime field. Two facts justify their use in \cref{lem:transfer}: their rational coefficients are integral at the chosen prime, and reduction preserves their total degrees.

\begin{theorem}[Rank-nine classification]\label{thm:r9}
Let $B\in\Q^9$ and $C\in\Q$. Then $F_{9,B,C}$ is a modular function if and only if $(B,C)$ is one of the fifteen pairs of \cref{conj:odd} with $m=4$.
\end{theorem}

\begin{proof}
Sufficiency follows from \cref{thm:basic,thm:k2,thm:k3to6}. For necessity, put
\[
t=1-\cos(\pi/11),\qquad f(t)=32t^5-144t^4+224t^3-140t^2+30t-1=0.
\]
The polynomial $f$ is irreducible, and the relevant embedding is its smallest real root, which lies between $1/25$ and $41/1000$. The exact saddle coordinates and covariance matrix lie in $K=\Q(t)$; their positivity, the saddle equations and the covariance inverse are checked over $K$.

Expand the three polynomials $c_1-\gamma_0$, $\kappa_2$ and $\kappa_3$ in the basis $1,t,t^2,t^3,t^4$ of $K$. Take the four nonconstant components of the first and all five components of each of the other two. These fourteen rational polynomials are necessary conditions for modularity. Introduce
\[
z=b_9,\qquad u_i=b_{2i-1}+b_{2i},\qquad x_i=b_{2i}+2z\qquad(1\le i\le4),
\]
write $u=(u_1,\dots,u_4)$ and $x=(x_1,\dots,x_4)$, and let $I$ be the ideal they generate in $\Q[u,x,z]$.

Choose $p=1000033$. All coefficients of the exact saddle quantities $Q_i$, $\xi_i$ and $M_{ij}$, written in the power basis of $K$, are integral at $p$. The leading coefficient and the discriminant of $f$ are units in $\Z_{(p)}$, and so are the Bernoulli and factorial denominators needed through order $h^6$. The recurrences for $H_a$, $D_a$ and the Gaussian moments therefore show that all fourteen polynomials have coefficients in $\Z_{(p)}$, and their denominators can be cleared by integers prime to $p$.

Modulo $p$, the polynomial $f$ has five distinct roots. We evaluate the Gaussian recurrence at each root and recover the power-basis components by inverse Vandermonde interpolation; comparison with the characteristic-zero saddle data identifies each embedding. The resulting polynomials modulo $p$ have degrees
\[
2,2,2,2,\quad3,3,3,3,3,\quad4,4,4,4,4.
\]
By \cref{lem:degree} these are also upper bounds for the degrees in characteristic zero. Hence reduction preserves the degrees, and the computed highest-degree parts are the reductions of the highest-degree parts of the rational polynomials.

An exact $F_4$ calculation over $\F_p$ gives the reduced ideal
\begin{equation}\label{eq:r9ideal}
\bar I=\bigl(u_i,\ x_i^2-x_i\ (1\le i\le4),\ x_ix_j-x_i\ (1\le i<j\le4),\ 4z^3-z\bigr),
\end{equation}
whose quotient has dimension fifteen. The ideal generated by the highest-degree parts has Hilbert series
\[
1+9s+41s^2+124s^3+271s^4+431s^5+449s^6+106s^7,
\]
and in particular contains every monomial of degree eight. The affine computation uses exact sparse linear algebra in \textsf{msolve} \cite{msolve}, and the homogeneous one uses \textsc{Singular} \cite{Singular}. \Cref{lem:transfer} gives $\dim_\Q\Q[u,x,z]/I\le15$.

The fifteen modular pairs supplied by the companion theorems give fifteen distinct rational zeros of $I$, and their evaluation map is surjective onto $\Q^{15}$. Hence the quotient has dimension exactly fifteen, and $I$ is the radical ideal of these points; equivalently, $I$ is generated over $\Q$ by the polynomials displayed in \eqref{eq:r9ideal}. Thus $u_i=0$, the $x_i$ form an ordered Boolean chain, and $z\in\{-\frac12,0,\frac12\}$. The first-order equation, checked exactly over $K$ at all fifteen points, gives
\[
C=-\frac{57}{176}+\frac{k(10-k)}{22}-kz+\frac{17}4z^2,\qquad k=x_1+x_2+x_3+x_4.
\]
This completes the proof of necessity.
\end{proof}

\section{Theta formulas and the basic companions}\label{sec:theta}

\subsection{The simultaneous shift}
Let $m\ge1$, $r=2m+1$, $a=(a_1,\dots,a_m)\in\Q^m$ and $c\in\Q$. Define
\[
B(a,c)=(a_1,-a_1,\dots,a_m,-a_m,c)^\T,\qquad S_m(a,c;q)=F_{2m+1,B(a,c),0}(q).
\]
The following consequence of the rank reduction of Shi and Wang explains a relation visible in the rank-three and rank-five tables. In rank five it specializes to \cite[(5.45)]{SW26} after replacing $q$ by $q^4$.

\begin{lemma}[Simultaneous shift]\label{lem:simshift}
For every integer $h$,
\begin{equation}\label{eq:simshift}
S_m(a+2h\one,c+h;q)=q^{-h\sum_{j=1}^ma_j-mh^2}\,\frac{(-q^{c+h+1/2};q)_\infty}{(-q^{c+1/2};q)_\infty}\,S_m(a,c;q).
\end{equation}
\end{lemma}

\begin{proof}
For $d\in\frac12\Z$, put
\[
\Theta_a(d;q)=\sum_{\ell\in\Z}q^{(\ell-d)^2-a\ell}.
\]
The change of summation variable $\ell\mapsto\ell+h$ gives
\[
\Theta_{a+2h}(d;q)=q^{-ha-h^2-2hd}\,\Theta_a(d;q).
\]
By \cite[Theorem~1.2, (1.11)]{SW26}, setting $n_0=0$ yields
\begin{multline*}
S_m(a,c;q)=\frac{(-q^{c+1/2};q)_\infty}{(q;q)_\infty^m}\sum_{n_1,\dots,n_m\ge0}\frac{q^{\frac14\sum_{j=1}^m(n_j-n_{j-1})^2}\,q^{\sum_{j=1}^{m-1}(a_{j+1}-a_j)n_j+(c-a_m)n_m}}{\prod_{j=1}^m(q;q)_{n_j}}\\
\times\prod_{j=1}^m\Theta_{a_j}\bigl((n_{j-1}-n_j)/2;q\bigr).
\end{multline*}
Under the simultaneous shift, the second exponential gains $q^{-hn_m}$. The product of theta factors gains $q^{-h\sum_ja_j-mh^2+hn_m}$, since $\sum_j(n_{j-1}-n_j)=-n_m$. These factors combine as asserted. All changes of variables are justified by absolute convergence.
\end{proof}

In particular, let $x_j\in\{0,1\}$ and $k=\sum_jx_j$, and define
\[
B_\pm=(\pm1-x_1,\ x_1\mp1,\ \dots,\ \pm1-x_m,\ x_m\mp1,\ \pm\tfrac12)^\T.
\]
Taking $h=-1$ in \cref{lem:simshift} gives the exact identity
\begin{equation}\label{eq:halfshift}
F_{r,B_-,C+k}(q)=2F_{r,B_+,C}(q)\qquad(C\in\Q).
\end{equation}
Thus modularity in the two half-shift families is equivalent. No ordering of the $x_j$ is needed for this identity. As one consequence, the all-rank modularity of the last-coordinate half-shift \cite{CTW26} also gives the modularity, for every odd $r=2m+1$, of
\[
F_{r,(-2,2,\dots,-2,2,-1/2),\,m+\frac{r(4-r)}{24(r+2)}}.
\]

\subsection{The proposed Wronskians}
For the remainder of the paper, let $m\ge1$, $r=2m+1$, $L=2m+3$ and
\[
\nu=\frac{r(r-1)}2=m(2m+1).
\]
For $0\le k\le m$, take $x_1=\cdots=x_{m-k}=0$ and $x_{m-k+1}=\cdots=x_m=1$. Write $B(k,z)$ and $C(k,z)$ for the vector and scalar of \cref{conj:odd}, and put
\[
S^{(z)}_{m,k}(q)=F_{r,B(k,z),0}(q).
\]
It suffices to consider $z=0,\frac12$, since \eqref{eq:halfshift} determines the series with $z=-\frac12$.

Define increasing tuples by deleting the indicated entry:
\begin{align*}
\rho^{(0)}_k&=(1,2,\dots,m+1)\setminus\{m+1-k\}, & \rho^{(0)}_*&=(1,2,\dots,m),\\
\rho^{(1/2)}_k&=(\tfrac12,\tfrac32,\dots,m+\tfrac12)\setminus\{k+\tfrac12\}, & \rho^{(1/2)}_*&=(\tfrac12,\tfrac32,\dots,m-\tfrac12).
\end{align*}
Deletion preserves the order of the remaining entries. For $v=(v_1,\dots,v_m)$, define
\[
\chi(v)=\prod_{i=1}^mv_i\prod_{1\le i<j\le m}(v_i^2-v_j^2),\qquad \|v\|^2=\sum_{i=1}^mv_i^2.
\]

\begin{conjecture}[Companion Wronskian identity]\label{conj:wron}
For $z\in\{0,\frac12\}$ and $0\le k\le m$, writing $\rho=\rho^{(z)}_k$, one has
\begin{equation}\label{eq:wron}
S^{(z)}_{m,k}(q)=\frac{(-q^{1/2+z};q)_\infty^r}{(q;q)_\infty^\nu}\sum_{v\in\rho+L\Z^m}\frac{\chi(v)}{\chi(\rho^{(z)}_*)}\,q^{\frac{\|v\|^2-\|\rho\|^2}{2L}}.
\end{equation}
\end{conjecture}

The cases $(k,z)=(0,0)$ and $(m,\frac12)$ are the principal and last-coordinate half-shift formulas of Chern, Tran and Wakhare \cite[Theorem~1.4]{CTW26}, which are based on \cite{BW15}. \Cref{thm:basic} proves the cases $k\in\{0,1,m\}$, \cref{thm:k2} adds $k=2$, and \cref{thm:k3to6} adds $k=3,4,5,6$.

\begin{proposition}[Modularity of the proposed right-hand sides]\label{prop:rhs-modular}
Let $R^{(z)}_{m,k}$ denote the right-hand side of \eqref{eq:wron}. Then $q^{C(k,z)}R^{(z)}_{m,k}(q)$ is a modular function on a congruence subgroup. Consequently, \cref{conj:wron} implies the sufficiency direction of \cref{conj:odd}.
\end{proposition}

\begin{proof}
For rational $a$, define the unary theta series
\[
\theta_{a,L}(q)=\sum_{v\in a+L\Z}v\,q^{v^2/(2L)},\qquad \mathcal D=q\frac{d}{dq}.
\]
Expanding the determinant and using the Vandermonde identity gives
\[
\det\bigl(\mathcal D^{i-1}\theta_{\rho_j,L}\bigr)_{i,j=1}^m=\frac{(-1)^{\binom m2}}{(2L)^{\binom m2}}\sum_{v\in\rho+L\Z^m}\chi(v)\,q^{\|v\|^2/(2L)}.
\]
The theta series $\theta_{\rho_j,L}$ have weight $3/2$ on a common congruence subgroup. Their Wronskian has weight $3m/2+m(m-1)=\nu/2$: the lower-order derivative terms in its transformation law form a triangular matrix and do not alter the determinant. This is the Wronskian argument used in \cite{MW24,CTW26}.

Put
\[
E_0(\tau)=\frac{\eta(\tau)^2}{\eta(\tau/2)\eta(2\tau)},\qquad E_{1/2}(\tau)=\frac{\eta(2\tau)}{\eta(\tau)}.
\]
The product identities
\[
(-q^{1/2};q)_\infty=q^{1/48}E_0(\tau),\qquad (-q;q)_\infty=q^{-1/24}E_{1/2}(\tau)
\]
and direct evaluation of the sums of squares give
\[
C(k,0)=\frac{\|\rho^{(0)}_k\|^2}{2L}-\frac{r^2}{48},\qquad
C(k,\tfrac12)=\frac{\|\rho^{(1/2)}_k\|^2}{2L}+\frac{r(3-r)}{48}.
\]
It follows that
\[
q^{C(k,z)}R^{(z)}_{m,k}(q)=\frac{E_z(\tau)^r}{\eta(\tau)^\nu\,\chi(\rho^{(z)}_*)}\sum_{v\in\rho^{(z)}_k+L\Z^m}\chi(v)\,q^{\|v\|^2/(2L)}.
\]
The weights cancel. Theta series and eta quotients have finite multiplier systems on congruence subgroups and are meromorphic at the cusps, so restriction to a common congruence subgroup gives the assertion.
\end{proof}

\subsection{Leading coefficients}
The constants in \eqref{eq:wron} are forced by the Nahm sums; they are not fitted separately in each rank.

\begin{proposition}[Leading coefficients]\label{prop:leading}
The series $S^{(0)}_{m,k}$ and $S^{(1/2)}_{m,k}$ have no negative powers of $q$. Their constant terms are, respectively,
\[
d^{(0)}_{m,k}=\binom{2m}{k}-\binom{2m}{k-2},\qquad d^{(1/2)}_{m,k}=\binom{2m+1}{m-k}.
\]
These equal $\chi(\rho^{(z)}_k)/\chi(\rho^{(z)}_*)$, the constant terms of the proposed right-hand sides.
\end{proposition}

\begin{proof}
Put $n_0=0$ and $a_i=n_i-n_{i-1}$, so that $n_j=\sum_{i=1}^ja_i\ge0$. If $E$ denotes the exponent in the unnormalized Nahm summand, then
\[
E=\frac12\sum_{i=1}^ra_i^2+zn_r-\sum_{j=1}^m(2z-x_j)a_{2j}.
\]
For $z=0$ this can be written as
\[
E=\frac{n_r}2+\sum_{i\notin I}\frac{a_i(a_i-1)}2+\sum_{i\in I}\frac{a_i(a_i+1)}2,\qquad I=\{2m-2k+2,\,2m-2k+4,\,\dots,\,2m\}.
\]
For $z=\frac12$ it becomes
\[
E=\sum_{i\in J}\frac{a_i(a_i-1)}2+\sum_{i\notin J}\frac{a_i(a_i+1)}2,\qquad J=\{2,4,\dots,2m-2k\}.
\]
All terms are nonnegative. Since each denominator has constant term one, the constant terms count the paths for which every displayed term vanishes.

We evaluate these path counts. On the basis $e_h$ indexed by integers $h\ge0$, let $Ue_h=e_{h+1}$, let $Ve_0=0$ and $Ve_h=e_{h-1}$ for $h\ge1$. Put $U_+=\Id+U$, $U_-=\Id+V$ and $J_0=U_-U_+=2\Id+U+V$, where $\Id$ is the identity operator. Let $[e_h]f$ denote the coefficient of $e_h$ in $f$, and let $\lambda(f)$ be the sum of the coefficients of $f$. The two zero-energy path counts are
\[
[e_0]\,J_0^kU_+^{2(m-k)}e_0\qquad\text{and}\qquad \lambda\bigl(U_-^{2k}J_0^{m-k}e_0\bigr).
\]
Reflection at height $-1$ gives
\[
[e_a]\,J_0^je_0=\Delta_j(a):=\binom{2j}{j-a}-\binom{2j}{j-a-2}.
\]
Here and below $\binom Nb=0$ unless $0\le b\le N$. By the symmetry of $J_0$ and Vandermonde convolution, the first count is
\[
\sum_{a\ge0}\binom{2(m-k)}a\Delta_k(a)=\binom{2m}k-\binom{2m}{k-2}.
\]
Writing $s=m-k$, the second count is
\[
\sum_{a\ge0}\Delta_s(a)\sum_{d=0}^a\binom{2k}d=\sum_{d\ge0}\binom{2k}d\left(\binom{2s}{s-d}+\binom{2s}{s-d-1}\right)=\binom{2m+1}s.
\]
Finally, direct cancellation in the finite products defining $\chi$ gives the same binomial formulas for the two ratios in the statement. Every coordinate of $\rho^{(z)}_k$ lies strictly between $0$ and $L/2$, so this tuple is the unique point of its lattice coset with minimal norm. Thus its contribution is exactly the constant term of the right-hand side of \eqref{eq:wron}.
\end{proof}

\subsection{The basic companion theorem}

\begin{theorem}[Basic companions]\label{thm:basic}
For every $m\ge1$, identity \eqref{eq:wron} holds for $k\in\{0,1,m\}$ and $z\in\{0,\frac12\}$. Consequently, $F_{2m+1,B(k,z),C(k,z)}$ is a modular function for
\[
k\in\{0,1,m\},\qquad z\in\{-\tfrac12,0,\tfrac12\}.
\]
For $m\ge2$ these are nine distinct pairs.
\end{theorem}

The proof uses two identities of Bartlett and Warnaar \cite[Theorems~3.7 and 5.2]{BW15}; we give the intervening determinant calculation and the specialization factors explicitly.

\subsection{Hall--Littlewood boundary formulas}
Throughout the boundary calculations, $w=q^{1/2}$, $L=2m+3$ and $\nu=m(2m+1)$. From here to the end of \cref{sec:cert}, $x=(x_1,\dots,x_m)$ denotes a vector of Hall--Littlewood variables; the Boolean coordinates of \cref{conj:odd} do not appear in these sections. A partition $\lambda$ is a finite weakly decreasing sequence of positive integers; $|\lambda|$ is its sum, $\ell(\lambda)$ its length, and $m_j(\lambda)$ the multiplicity of $j$. Define
\begin{equation}\label{eq:Hdef}
\HL_m(t;x;q)=\sum_{\lambda:\ \lambda_1\le2}q^{|\lambda|/2}\,t^{m_1(\lambda)}\,P'_\lambda(x_1,x_1^{-1},\dots,x_m,x_m^{-1};q),
\end{equation}
including the empty partition, whose summand is $1$. Here $P'_\lambda(X;q)=P_\lambda[X/(1-q);q]$ is the modified Hall--Littlewood symmetric function: $X/(1-q)$ replaces each letter $X_i$ by $X_i,X_iq,X_iq^2,\dots$. The ordinary Hall--Littlewood functions are determined, for $N\ge\ell(\lambda)$, by
\[
P_\lambda(X_1,\dots,X_N;q)=\frac1{v_\lambda(q)}\sum_{\sigma\in\mathfrak S_N}\sigma\Bigl(X_1^{\lambda_1}\cdots X_N^{\lambda_N}\prod_{i<j}\frac{X_i-qX_j}{X_i-X_j}\Bigr),\qquad v_\lambda(q)=\prod_{j\ge0}\frac{(q;q)_{m_j(\lambda)}}{(1-q)^{m_j(\lambda)}}.
\]
Zero parts are appended to length $N$, so $m_0(\lambda)=N-\ell(\lambda)$, and $\mathfrak S_N$ acts by permuting the variables. Stable specialization defines the symmetric function used in \eqref{eq:Hdef}.

For generic $x=(x_1,\dots,x_m)$ put
\begin{align*}
D_m(x;q)&=(q;q)_\infty^m\prod_{i=1}^m(qx_i^2,qx_i^{-2};q)_\infty\prod_{i<j}(qx_ix_j,qx_i/x_j,qx_j/x_i,q/(x_ix_j);q)_\infty,\\
P_t(x)&=\frac{\prod_i(-wtx_i,-wt/x_i;q)_\infty}{D_m(x;q)},\\
\delta_C(x)&=\det\bigl(x_i^j-x_i^{-j}\bigr)_{i,j=1}^m,\qquad \delta_B(x)=\det\bigl(x_i^{j-1/2}-x_i^{-j+1/2}\bigr)_{i,j=1}^m.
\end{align*}
For the rational functions and half-integral Laurent polynomials below, at values without poles on $xq^\Z$, set
\[
\mu_x[f]=\sum_{\ell\in\Z}q^{L\ell^2/2}x^{L\ell}f(xq^\ell),\qquad \Phi_a(x)=\mu_x[y^a-y^{-a}],
\]
and
\[
\Psi_j(x)=\mu_x\left[\frac{y^j-y^{-j}}{1+w^2+w(y+y^{-1})}\right].
\]
The letter $y$ denotes the argument of $f$ inside $\mu_x$. Powers are taken with a fixed branch of the logarithm on a neighbourhood of the relevant points: the evaluations below are first made for $0<q<1$ and positive variables, and then continued analytically. Gaussian decay gives normal convergence away from the poles.

\begin{lemma}[Boundary specialization]\label{lem:boundary}
For rational $v,c$, at points where the expressions are defined,
\begin{equation}\label{eq:boundary}
S_m(-2v\one,c;q)=(-q^{c+1/2};q)_\infty\,\HL_m(q^{c+v};q^v\one;q).
\end{equation}
For generic $x$ one has
\begin{align}
\HL_m(1;x;q)&=P_1(x)\frac{\det(\Phi_j(x_i))_{i,j=1}^m}{\delta_C(x)},\label{eq:H1}\\
\HL_m(q;x;q)&=P_1(x)\frac{\det(\Psi_j(x_i))_{i,j=1}^m}{\delta_C(x)},\label{eq:Hq}\\
\HL_m(w;x;q)&=P_w(x)\frac{\det(\Phi_{j-1/2}(x_i))_{i,j=1}^m}{\delta_B(x)}.\label{eq:Hw}
\end{align}
Specializations at coincident variables mean the removable limits.
\end{lemma}

\begin{proof}
Sum over the last coordinate of the Nahm sum using Euler's identity
\[
\sum_{n\ge0}\frac{q^{n(n-1)/2}z^n}{(q;q)_n}=(-z;q)_\infty.
\]
The resulting factor is $(-q^{c+1/2-h};q)_\infty$, where $h$ is the remaining last coordinate. Extract $(-q^{c+1/2};q)_\infty$ and reverse the order of the remaining $2m$ coordinates. In the notation of \cite[(3.21)]{BW15}, the series becomes $F_{1,2m}(q^v,0,q^{c+v};q)$. Their Theorem~3.7 holds for these parameters and identifies this expression with the series on the right of \eqref{eq:boundary}.

For the determinant formulas, take the level parameter equal to $1$ in \cite[(5.6a)]{BW15}, let their parameter $b$ tend to infinity, and set their parameter $c$ equal to $-w/t$. Their equation (5.7) then gives \eqref{eq:Hdef}. After rewriting the type-$C$ denominator as $\delta_C$, the summand has one-variable multiplier
\[
q^{L\ell^2/2}x^{L\ell}t^\ell\,\frac{(-wx/t;q)_\ell}{(-wtx;q)_\ell},
\]
where $(a;q)_\ell=(a;q)_\infty/(aq^\ell;q)_\infty$ for $\ell\in\Z$. The quotient is $1$ at $t=1$, which proves \eqref{eq:H1}. At $t=q$, combine its two telescoping factors with $P_q/P_1$; the resulting multiplier is
\[
\frac{y/w}{(1+y/w)(1+wy)}=\frac1{1+w^2+w(y+y^{-1})},\qquad y=xq^\ell.
\]
This proves \eqref{eq:Hq}. At $t=w$ the quotient is $(1+x)/(1+xq^\ell)$. Using
\[
\delta_C(y)=\prod_i\bigl(y_i^{1/2}+y_i^{-1/2}\bigr)\,\delta_B(y)
\]
to cancel this quotient and $w^\ell$ row by row gives \eqref{eq:Hw}. These calculations hold initially where the cited bilateral series converge without poles; analytic continuation gives the stated specializations. The limits needed below are computed explicitly.
\end{proof}

\subsection{A finite tridiagonal calculation}
For an increasing tuple $a=(a_1,\dots,a_m)$ define
\[
J_a(x)=P_1(x)\frac{\det(\Phi_{a_j}(x_i))_{i,j=1}^m}{\delta_C(x)},\qquad J_k(x)=J_{\rho^{(0)}_k}(x).
\]

\begin{lemma}\label{lem:tridiag}
For generic $x$,
\begin{equation}\label{eq:tridiag}
\HL_m(1;x;q)-\HL_m(q;x;q)=wJ_1(x).
\end{equation}
\end{lemma}

\begin{proof}
Put $d=m+1$, so that $L=2d+1$. Reindexing the bilateral sum gives
\[
\mu_x[f(y)]=q^{L/2}\mu_x[y^Lf(qy)].
\]
Together with $q=w^2$ this implies
\[
\Psi_{d+1}+w\Psi_d=\mu_x\left[\frac{y^{d+1}}{1+wy}-\frac{y^{-d-1}}{1+w/y}\right]=0,
\]
since the displayed reindexing transforms the second rational term into the first. Multiplying by the denominator in the definition of $\Psi_j$ also gives
\[
\Phi_j=(1+w^2)\Psi_j+w(\Psi_{j-1}+\Psi_{j+1}),\qquad \Psi_0=0.
\]
Thus $(\Phi_1,\dots,\Phi_d)^\T=U_d(\Psi_1,\dots,\Psi_d)^\T$, where $U_d$ is tridiagonal with entries $w$ above and below the diagonal and diagonal $1+w^2,\dots,1+w^2,1$. Direct multiplication gives
\[
\det U_d=1,\qquad (U_d^{-1})_{ij}=(-w)^{|i-j|}\sum_{s=0}^{\min(i,j)-1}w^{2s}.
\]
By the complementary-minor formula, the minor of $U_d^{-1}$ formed by its first $m$ rows with column $t$ deleted is $(-1)^{d+t}(U_d)_{t,d}$. It equals $1$ for $t=d$, $-w$ for $t=d-1$, and $0$ otherwise. The Cauchy--Binet formula therefore yields
\[
\det(\Psi_j(x_i))_{i,j=1}^m=\det(\Phi_1,\dots,\Phi_m)-w\det(\Phi_1,\dots,\Phi_{m-1},\Phi_{m+1}),
\]
where each determinant evaluates its columns at $x_1,\dots,x_m$. Equations \eqref{eq:H1} and \eqref{eq:Hq} now prove the lemma.
\end{proof}

\subsection{Confluent limits and the affine translation}
For an increasing tuple $a$ put
\[
N_a(q)=\sum_{v\in a+L\Z^m}\chi(v)\,q^{(\|v\|^2-\|a\|^2)/(2L)},\qquad A_0=\frac{(-w;q)_\infty^{2m}}{(q;q)_\infty^\nu},\qquad A_{1/2}=\frac{(-q;q)_\infty^{2m}}{(q;q)_\infty^\nu}.
\]

\begin{lemma}\label{lem:confluent}
For every $0\le k\le m$,
\begin{align}
J_k(\one)&=A_0\,\frac{N_{\rho^{(0)}_k}}{\chi(\rho^{(0)}_*)},\label{eq:conf0}\\
J_k(w\one)&=q^{-k/2}A_{1/2}\,\frac{N_{\rho^{(1/2)}_k}}{\chi(\rho^{(1/2)}_*)}.\label{eq:confhalf}
\end{align}
Moreover,
\begin{align}
\HL_m(w;\one;q)&=A_{1/2}\,\frac{N_{\rho^{(1/2)}_m}}{\chi(\rho^{(1/2)}_*)},\label{eq:Hw1}\\
\HL_m(w;w\one;q)&=A_0\,\frac{N_{\rho^{(0)}_m}}{\chi(\rho^{(0)}_*)}.\label{eq:Hww}
\end{align}
\end{lemma}

\begin{proof}
Set $x_i=e^{\varepsilon t_i}$ with distinct generic $t_i$, and expand each column in odd powers of $\varepsilon t_i$. Cancellation of the leading Vandermonde determinants gives
\[
\lim_{x\to\one}\frac{\det(\Phi_{a_j}(x_i))}{\delta_C(x)}=\frac{N_a}{\chi(\rho^{(0)}_*)}.
\]
The same calculation with $\delta_B$ gives the denominator $\chi(\rho^{(1/2)}_*)$. Normal convergence justifies the Taylor expansions. Evaluating the nonsingular infinite products proves \eqref{eq:conf0} and \eqref{eq:Hw1}.

Reindexing the theta series gives the exact affine translation
\begin{equation}\label{eq:affine}
\Phi_a(wu)=-q^{-a/2}u^{-L/2}\,\Phi_{L/2-a}(u).
\end{equation}
Write $h=\binom m2$ and $s=m+h$. Product formulas for the two finite determinants give, with $Z(u)=\prod_i(1-u_i^{-2})\prod_{i<j}(1-(u_iu_j)^{-1})$,
\begin{align*}
\lim_{u\to\one}\frac{D_m(wu;q)}{Z(u)}&=\frac{(q;q)_\infty^\nu}{(1-q)^s},\\
\lim_{u\to\one}\frac{\delta_C(wu)Z(u)}{\delta_C(u)}&=(w-w^{-1})^s,\\
\lim_{u\to\one}\frac{\delta_B(wu)Z(u)}{\delta_C(u)}&=(w^{1/2}-w^{-1/2})^m(w-w^{-1})^h.
\end{align*}
For example, $\delta_C(u)=\prod_i(u_i-u_i^{-1})\prod_{i<j}(u_j-u_i)(1-(u_iu_j)^{-1})$, which makes the cancellations explicit; the analogous formula for $\delta_B$ replaces each first factor by $u_i^{1/2}-u_i^{-1/2}$. Consequently,
\begin{align*}
\lim_{u\to\one}\frac{P_1(wu)\,\delta_C(u)}{\delta_C(wu)}&=(-1)^s2^mq^{s/2}A_{1/2},\\
\lim_{u\to\one}\frac{P_w(wu)\,\delta_C(u)}{\delta_B(wu)}&=(-1)^sq^{m/4+h/2}A_0.
\end{align*}
In \eqref{eq:affine}, reversing the order of the transformed columns contributes $(-1)^h$, and the $m$ minus signs contribute $(-1)^m$; together these cancel $(-1)^s$. For $a=\rho^{(0)}_k$ the transformed increasing tuple is $\rho^{(1/2)}_k$, and $\sum_ja_j=s+k$. Since $\chi(\rho^{(1/2)}_*)=2^{-m}\chi(\rho^{(0)}_*)$, the remaining scalar is exactly $q^{-k/2}$, as in \eqref{eq:confhalf}. For $a=\rho^{(1/2)}_*$ the transformed tuple is $(2,3,\dots,m+1)=\rho^{(0)}_m$, and $\sum_ja_j=m^2/2$; the scalar cancels because $m/4+h/2=m^2/4$. Applying \eqref{eq:Hw} proves \eqref{eq:Hww}. These product cancellations also show that all the limits used exist.
\end{proof}

\subsection{Contiguity and completion of the proof}
For brevity write $F_B=F_{2m+1,B,0}$. Subtracting adjacent summands in the defining series gives
\begin{equation}\label{eq:contig}
F_B-F_{B+e_i}=q^{(T_{2m+1})_{ii}/2+b_i}\,F_{B+T_{2m+1}e_i}.
\end{equation}
Take $B=B(a\one,c)$ and $E=e_{2m}-e_{2m-1}$. Apply \eqref{eq:contig} first at the last coordinate and then at coordinate $2m$ of the shifted vector. Since $T_{2m+1}(e_{2m}+e_{2m+1})=E$, the result is
\[
F_{B+E}=q^a\Bigl(q^{-c-1/2}F_B-(1+q^{-c-1/2})F_{B+e_{2m+1}}\Bigr).
\]
With $a=-2v$, $c=-v$ and \cref{lem:boundary}, this becomes
\begin{equation}\label{eq:contig2}
F_{B+E}=q^{-v-1/2}(-q^{1/2-v};q)_\infty\bigl(\HL_m(1;q^v\one;q)-\HL_m(q;q^v\one;q)\bigr).
\end{equation}

\begin{proof}[Proof of \cref{thm:basic}]
For $z=0$, \cref{lem:boundary} and \eqref{eq:contig2} at $v=0$ give
\[
S^{(0)}_{m,0}=(-w;q)_\infty J_0(\one),\qquad S^{(0)}_{m,1}=(-w;q)_\infty J_1(\one),
\]
the second via \cref{lem:tridiag}. For $k=m$, use \eqref{eq:boundary} with $v=\frac12$, $c=0$:
\[
S^{(0)}_{m,m}=(-w;q)_\infty\HL_m(w;w\one;q).
\]
Equations \eqref{eq:conf0} and \eqref{eq:Hww} prove the three identities with $z=0$.

At $v=\frac12$, $c=-\frac12$, identity \eqref{eq:boundary} concerns the negative half-shift with $k=0$. Using \eqref{eq:halfshift} and $(-1;q)_\infty=2(-q;q)_\infty$ we obtain
\[
S^{(1/2)}_{m,0}=(-q;q)_\infty J_0(w\one).
\]
Similarly, \eqref{eq:contig2} at $v=\frac12$ concerns the negative half-shift with $k=1$. Since $S^{(-1/2)}_{m,1}=2q^{-1}S^{(1/2)}_{m,1}$, it gives
\[
S^{(1/2)}_{m,1}=w(-q;q)_\infty J_1(w\one).
\]
Finally, \eqref{eq:boundary} at $v=0$, $c=\frac12$ yields
\[
S^{(1/2)}_{m,m}=(-q;q)_\infty\HL_m(w;\one;q).
\]
Equations \eqref{eq:confhalf} and \eqref{eq:Hw1} give the half-shift identities. \Cref{prop:rhs-modular} supplies modularity with exactly the scalar $C(k,z)$, and \eqref{eq:halfshift} supplies the sector $z=-\frac12$. When $m=1$ some of the indices coincide, consistently with the six rank-three pairs.
\end{proof}

For later use, set
\[
B_v=B(-2v\one,-v)\qquad(v\in\Q).
\]
By \cref{lem:boundary}, every last-coordinate shift of $B_v$ is a Hall--Littlewood boundary specialization.

\subsection{A worked example}\label{subsec:example}
We follow one pair through the paper: rank five ($m=2$, $L=7$, $\nu=10$), index $k=1$ and sector $z=\frac12$. In \cref{conj:odd} this is $x=(0,1)$, so
\[
B(1,\tfrac12)=(1,-1,0,0,\tfrac12)^\T,\qquad C(1,\tfrac12)=-\frac{55}{336}+\frac{5}{14}-\frac12+\frac9{16}=\frac{43}{168}.
\]
In the coordinates of \cref{thm:r5} this point is $u=v=0$, $(x,y)=(0,1)$, $z=\frac12$, and \eqref{eq:r5C} gives the same scalar $(120-168+189-55)/336=43/168$; it is the sixth row of the rank-five table, whose modularity is \cite[(5.34)]{SW26}.

On the companion side, $\rho^{(1/2)}_1=(\frac12,\frac52)$ and $\rho^{(1/2)}_*=(\frac12,\frac32)$, so $\chi(\rho^{(1/2)}_1)=-\frac{15}2$ and $\chi(\rho^{(1/2)}_*)=-\frac32$. Identity \eqref{eq:wron} reads
\begin{align*}
F_{5,(1,-1,0,0,1/2),0}(q)&=\frac{(-q;q)_\infty^5}{(q;q)_\infty^{10}}\sum_{v\in(\frac12,\frac52)+7\Z^2}\frac{v_1v_2(v_1^2-v_2^2)}{-3/2}\,q^{(v_1^2+v_2^2-13/2)/14}\\
&=5+45q+200q^2+715q^3+\cdots,
\end{align*}
and the constant term $5=\binom{5}{1}$ agrees with \cref{prop:leading}. This case is covered by \cref{thm:basic}: \eqref{eq:contig2} at $v=\frac12$ expresses the negative half-shift as a difference of two Hall--Littlewood boundary specializations, \cref{lem:tridiag} turns that difference into the single theta minor $wJ_1$, and \eqref{eq:confhalf} evaluates its confluent limit at $x=w\one$. By \cref{prop:rhs-modular}, $q^{43/168}$ times the right-hand side is a modular function, in agreement with the necessity computation.

\section{Polynomial boundary calculus}\label{sec:boundary}

The determinant side of the companion problem has a uniform solution. We first expand all boundary functions in a fixed basis of theta minors, then invert the expansion and fold indices beyond its dimension. For integers $0\le j\le k$, write
\begin{equation}\label{eq:Adef}
A_{k,j}(q)=(-1)^jq^{\binom{k-j}2}\qbin{k+j}{2j}.
\end{equation}

\subsection{A uniform recurrence for boundary determinants}
Retain $d=m+1$, $L=2d+1$ and $w=q^{1/2}$. For $j\ge0$, put
\[
W_j(y)=\prod_{s=0}^{j-1}(1+wq^sy)(1+wq^s/y),\qquad \Psi^{[j]}_a(x)=\mu_x\left[\frac{y^a-y^{-a}}{W_j(y)}\right].
\]
The empty product gives $\Psi^{[0]}_a=\Phi_a$. Let $U_d(t)$ be the $d\times d$ tridiagonal matrix with off-diagonal entries $t$ and diagonal $1+t^2,\dots,1+t^2,1$. Define
\[
V_0=I_d,\qquad V_j=U_d(w)U_d(wq)\cdots U_d(wq^{j-1})\quad(j\ge1),
\]
where $I_d$ is the identity matrix.

\begin{lemma}[Boundary determinant expansion]\label{lem:bdexp}
For generic $x$ and every integer $j\ge0$,
\begin{equation}\label{eq:bdexp}
\HL_m(q^j;x;q)=\sum_{k=0}^m(-1)^k(V_j)_{d-k,d}\,J_k(x).
\end{equation}
In particular, for $m\ge2$,
\begin{equation}\label{eq:bdexp2}
\HL_m(q^2;x;q)-(1+q+q^2)\HL_m(q;x;q)+q\HL_m(1;x;q)=q^2J_2(x).
\end{equation}
\end{lemma}

\begin{proof}
In the bilateral formula used in the proof of \cref{lem:boundary}, take $t=q^j$. The finite Pochhammer quotient telescopes over the indices $-j,\dots,j-1$. Combined with $P_{q^j}/P_1$ it gives $1/W_j(y)$, where $y=xq^\ell$. Thus
\[
\HL_m(q^j;x;q)=P_1(x)\frac{\det(\Psi^{[j]}_a(x_i))_{i,a=1}^m}{\delta_C(x)}.
\]

Fix $j\ge1$ and set $t=wq^{j-1}$. Direct cancellation shows
\[
\Psi^{[j]}_{d+1}+t\Psi^{[j]}_d=\mu_x\left[\frac{y^{d+1}}{(1+ty)W_{j-1}(y)}-\frac{y^{-d-1}}{(1+t/y)W_{j-1}(y)}\right]=0.
\]
To verify the last equality, apply $\mu_x[f(y)]=q^{L/2}\mu_x[y^Lf(qy)]$ to the second term and use
\[
(1+t/(qy))\,W_{j-1}(qy)=\frac{(1+ty)\,W_{j-1}(y)}{wy}.
\]
Multiplication by the last factor of $W_j$ also gives
\[
\Psi^{[j-1]}_a=(1+t^2)\Psi^{[j]}_a+t\bigl(\Psi^{[j]}_{a-1}+\Psi^{[j]}_{a+1}\bigr),\qquad \Psi^{[j]}_0=0.
\]
The boundary relation therefore closes this recurrence at $a=d$:
\[
\bigl(\Psi^{[j]}_1,\dots,\Psi^{[j]}_d\bigr)^\T=V_j^{-1}(\Phi_1,\dots,\Phi_d)^\T.
\]
Each $U_d(t)$ has determinant $1$, and the complementary-minor formula and the Cauchy--Binet formula give \eqref{eq:bdexp}.

The last column of $V_1$ has entries $w,1$ in its last two rows. The last column of $V_2$ has entries
\[
q^2,\qquad w(1+q+q^2),\qquad 1+q^2
\]
in its last three rows and zeros elsewhere. Substitution in \eqref{eq:bdexp} proves \eqref{eq:bdexp2}.
\end{proof}

\subsection{The half-line kernel and Bailey inversion}

\begin{proposition}[Explicit boundary coefficients]\label{prop:coeffs}
Let $d=m+1$ and $0\le h\le j\le m$. Then
\begin{equation}\label{eq:coeffs}
(V_j)_{d-h,d}=q^{h^2/2}D_{j,h}(q),\qquad D_{j,h}(q)=\qbin{2j}{j-h}-q^{2h+1}\qbin{2j}{j-h-1}.
\end{equation}
Gaussian binomial coefficients outside their usual range are zero.
\end{proposition}

\begin{proof}
We first work on the half-line with coordinates $0,1,2,\dots$, where coordinate zero corresponds to the last row of the finite matrix. Let $\mathsf U(t)$ be tridiagonal, with off-diagonal entries $t$, diagonal entry $1$ at zero, and diagonal entries $1+t^2$ elsewhere. All matrix products below have finite bandwidth, so each entry is a finite sum. Put
\[
f_j(a)=q^{a^2/2}\qbin{2j}{j+a}\quad(a\in\Z),\qquad K_j(h,\ell)=f_j(h-\ell)-q^{h+1/2}f_j(h+\ell+1).
\]
The finite $q$-binomial theorem gives the Laurent polynomial identity
\[
\sum_{a\in\Z}f_j(a)\,y^a=\prod_{s=0}^{j-1}(1+wq^sy)(1+wq^s/y).
\]
Consequently, with $t_j=wq^j$,
\[
f_{j+1}(a)=(1+t_j^2)f_j(a)+t_j\bigl(f_j(a-1)+f_j(a+1)\bigr).
\]
The same recurrence holds for $K_j(h,\ell)$ in the second coordinate. The elementary ratio of consecutive Gaussian binomial coefficients gives
\[
(1-q^{j+h+1})f_j(h+1)=q^{h+1/2}(1-q^{j-h})f_j(h),
\]
also at the endpoints, where the appropriate terms vanish. Equivalently,
\[
K_j(h,-1)=-t_jK_j(h,0).
\]
This relation changes the diagonal coefficient at zero from $1+t_j^2$ to $1$. Since $K_0(h,\ell)=\delta_{h\ell}$ for $h,\ell\ge0$, induction proves
\[
K_j(h,\ell)=\bigl(\mathsf U(w)\,\mathsf U(wq)\cdots\mathsf U(wq^{j-1})\bigr)_{h,\ell}.
\]
For the column $\ell=0$ and $j\le m$, no matrix path reaches a coordinate above $m$, so this column agrees with that of the finite product $V_j$. Substituting $\ell=0$ gives \eqref{eq:coeffs}.
\end{proof}

\begin{lemma}[Uniform determinant inverse]\label{lem:inverse}
For generic $x$ and every $0\le k\le m$,
\begin{equation}\label{eq:inverse}
\sum_{j=0}^kA_{k,j}(q)\,\HL_m(q^j;x;q)=q^{k^2/2}J_k(x).
\end{equation}
\end{lemma}

\begin{proof}
By \cref{lem:bdexp} and \cref{prop:coeffs}, it suffices to show that
\begin{equation}\label{eq:inv-coeff}
\sum_{j=h}^kA_{k,j}(q)\,D_{j,h}(q)=(-1)^k\delta_{kh}\qquad(0\le h\le k).
\end{equation}
The case $h=k$ is immediate. Suppose $n=k-h\ge1$. Direct cancellation gives
\[
D_{j,h}(q)=\frac{(1-q^{2h+1})(q;q)_{2j}}{(q;q)_{j-h}(q;q)_{j+h+1}}.
\]
Define a polynomial of degree $n-1$ by
\[
P(T)=\prod_{a=2k-n+2}^{2k}(1-q^aT).
\]
With $\ell=k-j$, the left-hand side of \eqref{eq:inv-coeff} becomes
\[
\frac{(-1)^k(1-q^{2h+1})}{(q;q)_n}\sum_{\ell=0}^n(-1)^\ell q^{\binom\ell2}\qbin n\ell P(q^{-\ell}).
\]
For each monomial $T^a$ of $P$ one has $0\le a<n$, and the finite $q$-binomial theorem evaluates its contribution as $(q^{-a};q)_n=0$. This proves \eqref{eq:inv-coeff} and the lemma.
\end{proof}

The coefficient inversion \eqref{eq:inv-coeff} is a rescaled form of classical Bailey inversion. Indeed, put
\[
\beta_j=\frac{\HL_m(q^j;x;q)}{(q;q)_{2j}},\qquad \alpha_h=\frac{1-q^{2h+1}}{1-q}(-1)^hq^{h^2/2}J_h(x).
\]
For $0\le j\le m$, \cref{prop:coeffs} gives
\[
\beta_j=\sum_{h=0}^j\frac{\alpha_h}{(q;q)_{j-h}(q^2;q)_{j+h}},
\]
which is the defining Bailey relation relative to $q$; see \cite[Lemma~1]{MZ10} for the more general well-poised inversion. We include the direct proof above to make the normalization and the finite range explicit.

\subsection{The finite kernel and boundary powers beyond the rank}
The half-line kernel also gives an exact formula for the finite matrix with no restriction on the number of factors. This supplies boundary-shift relations even when the index exceeds the rank parameter.

\begin{proposition}[Finite boundary kernel]\label{prop:kernel}
Put $d=m+1$ and $L=2d+1$. For $j\ge0$ and $a\in\Z$ define
\[
f_j(a)=q^{a^2/2}\qbin{2j}{j+a},\qquad \gamma_i(h)=q^{i(h+1/2)+Li(i-1)/2}\quad(i\in\Z).
\]
For $0\le h,\ell<d$,
\begin{equation}\label{eq:kernel}
(V_j)_{d-h,d-\ell}=\sum_{i\in\Z}\gamma_i(h)\Bigl(f_j(h-\ell+iL)-q^{h+iL+1/2}f_j(h+\ell+1+iL)\Bigr).
\end{equation}
The sum is finite because $f_j(a)=0$ for $|a|>j$.
\end{proposition}

\begin{proof}
Write $K_j(a,\ell)=f_j(a-\ell)-q^{a+1/2}f_j(a+\ell+1)$, now allowing $a\in\Z$, and denote the right-hand side of \eqref{eq:kernel} by $\widehat K_j(h,\ell)$. The recurrence and the lower boundary calculation in the proof of \cref{prop:coeffs} remain valid for every integer $a$. Thus, with $t_j=wq^j$,
\[
\widehat K_j(h,-1)=-t_j\widehat K_j(h,0).
\]
At the upper boundary, the two arguments of $f_j$ in $K_j(h+iL,d)$ differ by $L$. Since
\[
\gamma_{i+1}(h)=q^{h+iL+1/2}\gamma_i(h),
\]
successive terms cancel, giving $\widehat K_j(h,d)=0$. The free three-term recurrence therefore becomes multiplication on the right by $U_d(t_j)$ in reversed coordinates: the lower boundary has diagonal entry $1$, and the upper boundary is truncated. Finally, $\widehat K_0(h,\ell)=\delta_{h\ell}$, since in the stated range only the term with $i=0$ and $h=\ell$ contributes. Induction on $j$ proves the formula.
\end{proof}

\begin{theorem}[Affine folding of the boundary inverse]\label{thm:folding}
For generic $x$ and every $k\ge0$, put
\[
\mathcal Q_{m,k}(x;q)=q^{-k^2/2}\sum_{j=0}^kA_{k,j}(q)\,\HL_m(q^j;x;q).
\]
Write $k=bL+s$ with $b\ge0$ and $0\le s<L$. Then
\begin{equation}\label{eq:folding}
\mathcal Q_{m,k}(x;q)=(-1)^bq^{b(s+1/2)+L\binom b2}\begin{cases}J_s(x),&0\le s<d,\\ 0,&s=d,\\ -J_{L-1-s}(x),&d<s<L.\end{cases}
\end{equation}
In particular,
\[
\mathcal Q_{m,m+1}=0,\qquad \mathcal Q_{m,k+L}=-q^{k+1/2}\mathcal Q_{m,k}\quad(k\ge0).
\]
\end{theorem}

\begin{proof}
For $a\ge0$, the coefficient calculation \eqref{eq:inv-coeff} gives
\[
\sum_{j=0}^kA_{k,j}(q)\,K_j(a,0)=(-1)^kq^{k^2/2}\delta_{ak};
\]
this includes $a>k$, since all summands then vanish. The symmetry $f_j(-a)=f_j(a)$ also gives
\[
K_j(-a-1,0)=-q^{-a-1/2}K_j(a,0).
\]
Substitute the finite kernel into \cref{lem:bdexp} and perform the sum over $j$. For fixed $h$, only image indices satisfying $h+iL=k$ or $h+iL=-k-1$ can survive.

If $s<d$, the first equation has the unique solution $h=s$, $i=b$, and the coefficient of $J_h$ is $(-1)^{k+h}\gamma_b(h)$; this is the first case of \eqref{eq:folding}. If $s>d$, the second equation has the unique solution $h=L-1-s$, $i=-b-1$. Its coefficient is $-(-1)^{k+h}\gamma_{-b-1}(h)q^{-k-1/2}$, which simplifies to the third case. For $s=d$, neither equation has a solution with $0\le h<d$. The stated quasi-periodicity follows directly.
\end{proof}

\begin{corollary}[A boundary-shift relation]\label{cor:bdrel}
Let $m\ge1$, $r=2m+1$ and $v\in\Q$. Then
\begin{equation}\label{eq:bdrel}
\sum_{j=0}^{m+1}A_{m+1,j}(q)\,(-q^{1/2-v};q)_j\,F_{r,B_v+je_r,0}(q)=0.
\end{equation}
\end{corollary}

\begin{proof}
By \cref{lem:boundary}, each term after the coefficient $A_{m+1,j}$ equals $(-q^{1/2-v};q)_\infty\HL_m(q^j;q^v\one;q)$. Apply \cref{thm:folding} with $k=m+1$.
\end{proof}

\subsection{A minimal basis for the boundary shifts}
The finite kernel gives reductions directly among boundary-shifted Nahm sums, without introducing companion sums. The coefficients are polynomials in $q$, despite the negative powers in the inverse formula.

\begin{theorem}[Boundary basis]\label{thm:basis}
For $j\ge0$ and $0\le\ell\le m$, define
\begin{equation}\label{eq:Cjl}
C_{j,\ell}(q)=\sum_{h=\ell}^m(-1)^hq^{-h^2/2}(V_j)_{d-h,d}\,A_{h,\ell}(q),\qquad d=m+1.
\end{equation}
Then $C_{j,\ell}\in\Z[q]$ and, for generic $x$,
\begin{equation}\label{eq:basis}
\HL_m(q^j;x;q)=\sum_{\ell=0}^mC_{j,\ell}(q)\,\HL_m(q^\ell;x;q).
\end{equation}
For each fixed $q$ with $0<|q|<1$, the $m+1$ functions $\HL_m(1;x;q),\dots,\HL_m(q^m;x;q)$ are linearly independent over $\C$ as meromorphic functions of $x$. Consequently these reductions are unique, and no nonzero relation with coefficients independent of $x$ involves only the first $m+1$ boundary shifts. The relation at $j=m+1$ is, up to a scalar factor,
\[
\sum_{j=0}^{m+1}A_{m+1,j}(q)\,\HL_m(q^j;x;q)=0.
\]
\end{theorem}

\begin{proof}
Substitute \cref{lem:inverse} into \eqref{eq:bdexp}; interchanging two finite sums gives \eqref{eq:Cjl} and \eqref{eq:basis}.

A matrix path contributing to $(V_j)_{d-h,d}$ has at least $h$ off-diagonal steps. Their time indices are distinct and nonnegative, and a step at time $s$ has weight $w^{2s+1}$, so their total exponent is at least $1+3+\cdots+(2h-1)=h^2$. Every diagonal contribution has even exponent in $w$, while the number of off-diagonal steps has the parity of $h$. Thus
\[
q^{-h^2/2}(V_j)_{d-h,d}\in\Z[q].
\]
Since $A_{h,\ell}\in\Z[q]$, the same holds for $C_{j,\ell}$.

To prove independence, fix $q$ and recall that
\[
\Phi_a(x)=\sum_{n\in\Z}q^{Ln^2/2}\bigl(q^{an}x^{Ln+a}-q^{-an}x^{Ln-a}\bigr),\qquad 1\le a\le d,\quad L=2d+1.
\]
The residue classes $\pm1,\dots,\pm d$ modulo $L$ are distinct, so the coefficient of $x^a$ in a linear combination of these Laurent series is exactly the coefficient of $\Phi_a$. Thus $\Phi_1,\dots,\Phi_d$ are linearly independent on $\C^\times$.

There exist $d$ evaluation points at which the matrix of values of $\Phi_1,\dots,\Phi_d$ is invertible. Evaluating the $d$ maximal minors on the $d$ subsets of $d-1=m$ of these points produces, up to signs, the cofactor matrix, which is also invertible. Hence these minors are linearly independent as functions of $m$ variables. Multiplication by the common nonzero meromorphic factor $P_1(x)/\delta_C(x)$ proves the independence of $J_0,\dots,J_m$.

For $0\le j\le m$, the coefficient matrix in \eqref{eq:bdexp} is triangular: the coefficient of $J_h$ vanishes for $h>j$, and the coefficient of $J_j$ is $(-1)^jq^{j^2/2}\ne0$. Therefore the first $m+1$ boundary functions are independent. The remaining assertions follow from this independence and \cref{thm:folding} with $k=m+1$.
\end{proof}

\begin{corollary}[Polynomial reduction of every boundary shift]\label{cor:reduction}
Let $m\ge1$, $r=2m+1$, $v\in\Q$, $b=q^{1/2-v}$ and $F_j=F_{r,B_v+je_r,0}$. For every integer $j\ge0$,
\begin{equation}\label{eq:reduction}
(-b;q)_jF_j=\sum_{\ell=0}^mC_{j,\ell}(q)\,(-b;q)_\ell F_\ell.
\end{equation}
In particular, $C_{j,\ell}=\delta_{j\ell}$ for $0\le j\le m$, and $C_{m+1,\ell}=(-1)^mA_{m+1,\ell}$.
\end{corollary}

\begin{proof}
\Cref{lem:boundary} gives $(-b;q)_jF_j=(-b;q)_\infty\HL_m(q^j;q^v\one;q)$. Apply \eqref{eq:basis} and take the removable limit at coincident variables. The last assertions follow from the uniqueness in \cref{thm:basis}.
\end{proof}

The minimality statement concerns free variables $x$. We do not claim independence after setting $x=q^v\one$, or after fixing one of the modular sectors.

\subsection{A polynomial recurrence}
For indeterminates $e,c$, define
\begin{equation}\label{eq:Rdef}
R_k(e;c)=\sum_{j=0}^kA_{k,j}(q)\,(-c;q)_j\,e^j.
\end{equation}
The specialization $c=w/u$ occurs in the companion certificates. The following recurrence holds before that specialization.

\begin{proposition}[Boundary polynomial recurrence]\label{prop:recurrence}
One has $R_0(e;c)=1$, $R_1(e;c)=1-(1+c)e$, and, for $k\ge2$,
\begin{multline}\label{eq:recurrence}
R_k(e;c)=(1+q)q^{k-1}R_{k-1}(e/q;c)-q^{2k-2}R_{k-2}(e/q^2;c)\\
-(1+c)q^{2k-2}e\,R_{k-1}(e/q^2;cq).
\end{multline}
The equality is an identity in $\Z[q,q^{-1},c,e]$, and its left-hand side belongs to $\Z[q,c,e]$.
\end{proposition}

\begin{proof}
Put $S_k(e)=\sum_{j=0}^kA_{k,j}(q)e^j$ and define
\[
\mathcal F_n(X,s)=\sum_{a=0}^{\lfloor n/2\rfloor}q^{\binom a2}\qbin{n-a}{a}s^aX^{n-2a}.
\]
These are Cigler's $q$-Fibonacci polynomials, in the normalization $\mathcal F_n(X,s)=F_{n+1}(X,s/q)$ of \cite[Theorem~2.1]{Cigler03}. The $q$-Pascal identity gives
\[
\mathcal F_n(X,s)=X\mathcal F_{n-1}(X,qs)+s\mathcal F_{n-2}(X,qs)\qquad(n\ge2).
\]
Eliminating the odd index by two applications of this recurrence yields
\[
\mathcal F_{2k}(X,s)=X^2\mathcal F_{2k-2}(X,q^2s)+(1+q)s\,\mathcal F_{2k-2}(X,qs)-q^2s^2\mathcal F_{2k-4}(X,q^2s).
\]
When $X^2=-e$, the change of index $a=k-j$ gives $\mathcal F_{2k}(X,1)=S_k(e)$, and homogeneity gives $\mathcal F_{2h}(X,q^a)=q^{ah}S_h(e/q^a)$. Consequently,
\[
S_k(e)=(1+q)q^{k-1}S_{k-1}(e/q)-q^{2k-2}S_{k-2}(e/q^2)-q^{2k-2}e\,S_{k-1}(e/q^2).
\]
Let $\mathcal L_c$ be the $\Z[q,q^{-1}]$-linear map determined by $\mathcal L_c(e^j)=(-c;q)_je^j$. It commutes with rescaling of $e$ and satisfies, for every polynomial $P(e)$,
\[
\mathcal L_c(eP(e))=(1+c)\,e\,\mathcal L_{cq}(P(e)).
\]
Applying this map to the recurrence for $S_k$ proves \eqref{eq:recurrence}. The initial values follow from \eqref{eq:Rdef}.
\end{proof}

The recurrence is an exact relation among the coefficient polynomials. It does not by itself prove a recurrence among the Nahm sums to which those polynomials are applied.

\section{Certified companion identities}\label{sec:cert}

\subsection{The second companion}

\begin{theorem}\label{thm:k2}
For every $m\ge2$, identity \eqref{eq:wron} holds for $k=2$ and $z=0,\frac12$. Consequently, the three pairs $(B(2,z),C(2,z))$ with $z=-\frac12,0,\frac12$ lie in $\M_{2m+1}$.
\end{theorem}

Let $v\in\Q$ and $u=q^v$, and set
\[
E_i=e_{2i}-e_{2i-1},\qquad G_{m,2}(v;q)=F_{2m+1,B_v+E_{m-1}+E_m,0}(q).
\]

\begin{lemma}\label{lem:k2red}
For every $m\ge2$,
\begin{multline}\label{eq:k2red}
G_{m,2}(v;q)=q^{-2v-2}(-q^{1/2-v};q)_\infty\Bigl(\HL_m(q^2;u\one;q)\\
-(1+q+q^2)\HL_m(q;u\one;q)+q\HL_m(1;u\one;q)\Bigr).
\end{multline}
\end{lemma}

\begin{proof}
Only the last five coordinates are used. Introduce indeterminates $a,b,c,d_0,e$ corresponding to $q^{n_{r-4}},\dots,q^{n_r}$, respectively, and put $M=abcd_0e$, $K=1+w/u$ and $R=Ke-1$. For a Laurent polynomial $P$, let $\sigma_bP$ be obtained by replacing $b$ by $qb$, and define $\sigma_c,\sigma_{d_0},\sigma_e$ similarly. Consider the four operators
\begin{align*}
\mathcal A_2P&=ac(1-b)P-u^2b^2\sigma_bP,\\
\mathcal A_3P&=bd_0(1-c)P-u^{-2}c^2\sigma_cP,\\
\mathcal A_4P&=ce(1-d_0)P-u^2d_0^2\sigma_{d_0}P,\\
\mathcal A_5P&=d_0(1-e)P-(wu)^{-1}e\,\sigma_eP,
\end{align*}
and the four Laurent polynomials
\begin{align*}
P_2&=d_0e\bigl(Ke^2-(K(1+w^4)+w^2)e+w^2\bigr)+\frac{w^3}ue^2R-w^3ud_0^2eR,\\
P_3&=w^3ua\,d_0eR,\\
P_4&=a\left(\frac{w^3}u\,beR-\frac{w^4}{u^2}\,ce-w^4Ke+w^3uKd_0(1-e)\right),\\
P_5&=ace\left(Ke-w^2(w^2+wu+1)-\frac{w^3}u\,c\right)+u^2b^2e\bigl(w^3uKd_0+w^2-Ke\bigr)+w^5u^3Ka\,d_0^2.
\end{align*}
Direct expansion gives the Laurent-polynomial identity
\begin{multline}\label{eq:k2cert}
\sum_{i=2}^5\mathcal A_iP_i=w^4u^2b^2d_0^2e\\
-M\Bigl((1+w/u)(1+w^3/u)e^2-(1+w^2+w^4)(1+w/u)e+w^2\Bigr).
\end{multline}
The four displayed polynomials constitute the complete certificate.

For a Laurent polynomial $P$ in these five variables, let
\[
\Lambda(P)=\sum_{n\in\Z_{\ge0}^r}\frac{q^{n^\T T_rn/2+B_v^\T n}}{\prod_{i=1}^r(q;q)_{n_i}}\,P(q^{n_{r-4}},\dots,q^{n_r}).
\]
The reindexing used in \eqref{eq:contig} gives $\Lambda(M^{-1}\mathcal A_iP)=0$ for $i=2,3,4,5$. For example, the shift in $b$ has factor $q^{1+2v}$ before conjugation by $M^{-1}$; the latter supplies $q^{-1}$, leaving the factor $u^2$ displayed in $\mathcal A_2$. At the last coordinate the factor is $q^{1/2-v-1}=(wu)^{-1}$. The factors $1-q^{n_i}$ remove the lower boundary terms, and positive definiteness justifies all reindexings, including those with Laurent-polynomial weights.

Apply $\Lambda(M^{-1}\,\cdot\,)$ to \eqref{eq:k2cert}. Writing $F_j=F_{r,B_v+je_r,0}$ gives
\[
q^{2v+2}G_{m,2}(v;q)=(1+q^{1/2-v})(1+q^{3/2-v})F_2-(1+q+q^2)(1+q^{1/2-v})F_1+qF_0.
\]
\Cref{lem:boundary} expresses $F_j$ as $(-q^{j+1/2-v};q)_\infty\HL_m(q^j;u\one;q)$. Extracting $(-q^{1/2-v};q)_\infty$ proves the assertion.
\end{proof}

\begin{proof}[Proof of \cref{thm:k2}]
\Cref{lem:bdexp,lem:k2red} give
\begin{equation}\label{eq:k2final}
G_{m,2}(v;q)=q^{-2v}(-q^{1/2-v};q)_\infty\,J_2(q^v\one).
\end{equation}
At $v=0$ this is $S^{(0)}_{m,2}=(-w;q)_\infty J_2(\one)$, and \eqref{eq:conf0} gives the claimed Wronskian formula. At $v=\frac12$ the left-hand side is $S^{(-1/2)}_{m,2}$. Using \eqref{eq:halfshift} and $(-1;q)_\infty=2(-q;q)_\infty$ in \eqref{eq:k2final} yields
\[
S^{(1/2)}_{m,2}=q(-q;q)_\infty J_2(w\one).
\]
The factor $q^{-1}$ in \eqref{eq:confhalf} with $k=2$ cancels $q$, which proves the positive half-shift formula. Modularity follows from \cref{prop:rhs-modular}, and the negative sector follows from \eqref{eq:halfshift}.
\end{proof}

\subsection{The third through sixth companions}
The third through sixth companions admit certificates of the same kind as \eqref{eq:k2cert}. Because of their size, their coefficients are recorded in the accompanying data (\cref{app:verification}).

\begin{theorem}[Third through sixth companions]\label{thm:k3to6}
Let $k\in\{3,4,5,6\}$ and $m\ge k$. Identity \eqref{eq:wron} holds for both $z=0$ and $z=\frac12$. Consequently, the three pairs $(B(k,z),C(k,z))$ with $z=-\frac12,0,\frac12$ lie in $\M_{2m+1}$.
\end{theorem}

Retain $u=q^v$, $B_v=B(-2v\one,-v)$ and $E_i=e_{2i}-e_{2i-1}$. Define
\[
G_{m,k}(v;q)=F_{2m+1,\,B_v+\sum_{i=m-k+1}^mE_i,\,0}(q),
\]
and abbreviate $R_k(e)=R_k(e;w/u)$ from \eqref{eq:Rdef}.

\begin{lemma}[Certified boundary reductions]\label{lem:certred}
Suppose either $k\in\{3,4\}$ and $v\in\Q$, or $k\in\{5,6\}$ and $v\in\{0,\frac12\}$. For every $m\ge k$,
\begin{equation}\label{eq:certred}
G_{m,k}(v;q)=q^{-kv-k^2/2}(-q^{1/2-v};q)_\infty\sum_{j=0}^kA_{k,j}(q)\,\HL_m(q^j;u\one;q).
\end{equation}
\end{lemma}

\begin{proof}
Only the last $2k+1$ coordinates occur. Write $a_i=q^{n_{r-2k+i}}$ for $0\le i\le2k$, and put $M_k=\prod_{i=0}^{2k}a_i$. In the algebraic calculation the $a_i$, $w$ and $u$ are independent indeterminates when $k=3,4$; for $k=5,6$ we specialize $u=1$ or $u=w$. Let $\sigma_i$ replace $a_i$ by $qa_i$, with $q=w^2$, and define
\begin{align*}
\mathcal T_iP&=a_{i-1}a_{i+1}(1-a_i)P-u^{2(-1)^{i+1}}a_i^2\sigma_iP\qquad(1\le i<2k),\\
\mathcal T_{2k}P&=a_{2k-1}(1-a_{2k})P-(wu)^{-1}a_{2k}\sigma_{2k}P.
\end{align*}
For $k=3,4$, the files \file{uniform_k3_certificate.json} and \file{uniform_k4_certificate.json} define Laurent polynomials $P_1,\dots,P_{2k}$ as finite coefficient lists: an entry with index $i$, exponent vector $\alpha$ and coefficient $c(w,u)$ contributes $c(w,u)\prod_{h=0}^{2k}a_h^{\alpha_h}$ to $P_i$. For $k=5$ the corresponding lists are \file{uniform_k5_zero_certificate.json} at $u=1$ and \file{uniform_k5_half_certificate.json} at $u=w$; their coefficients are Laurent polynomials in $w$ alone. For $k=6$, the files \file{uniform_k6_zero_certificate.json} and \file{uniform_k6_half_certificate.json} contain integer polynomials in $w$, and each also specifies, in its field \file{target_multiplier}, a nonzero polynomial $D_{6,v}(w)$. Set $D_{k,v}=1$ for $k=3,4,5$. The coefficient lists satisfy the cleared identity
\begin{equation}\label{eq:certid}
\sum_{i=1}^{2k}\mathcal T_iP_i=D_{k,v}(w)\Bigl(w^{k^2}u^ka_{2k}\prod_{h=0}^{k-1}a_{2h+1}^2-M_kR_k(a_{2k})\Bigr).
\end{equation}
Identity \eqref{eq:certid} is checked by exact expansion in the Laurent polynomial ring over $\Z$, with all remaining variables independent; for $k=6$ the expansion is carried out one tail monomial at a time, and $D_{6,v}\ne0$ is checked directly.

Define $\Lambda$ as in the proof of \cref{lem:k2red}, using these $2k+1$ tail coordinates and the base shift $B_v$. After conjugation by $M_k^{-1}$, the reindexing factors are $u^{2(-1)^{i+1}}$ at an interior coordinate and $(wu)^{-1}$ at the last coordinate. Hence $\Lambda(M_k^{-1}\mathcal T_iP_i)=0$. The factors $1-a_i$ remove the lower boundary terms, and positive definiteness permits the shifts in the absolutely convergent sums. Applying $\Lambda(M_k^{-1}\,\cdot\,)$ to \eqref{eq:certid} and cancelling $D_{k,v}(w)$ wherever it is nonzero gives
\[
q^{kv+k^2/2}G_{m,k}(v;q)=\sum_{j=0}^kA_{k,j}(q)\,(-q^{1/2-v};q)_j\,F_{r,B_v+je_r,0}(q).
\]
Both sides are holomorphic on the upper half-plane, so the identity theorem extends the equality across the zeros of $D_{k,v}(w)$. By \cref{lem:boundary},
\[
F_{r,B_v+je_r,0}(q)=(-q^{j+1/2-v};q)_\infty\HL_m(q^j;u\one;q).
\]
Combining the finite and infinite products gives \eqref{eq:certred}.
\end{proof}

\begin{proof}[Proof of \cref{thm:k3to6}]
Combining \cref{lem:certred,lem:inverse} yields
\[
G_{m,k}(v;q)=q^{-kv}(-q^{1/2-v};q)_\infty J_k(q^v\one).
\]
At $v=0$, the confluent limit \eqref{eq:conf0} gives \eqref{eq:wron} for $z=0$. At $v=\frac12$, the half-shift identity \eqref{eq:halfshift} and $(-1;q)_\infty=2(-q;q)_\infty$ give
\[
S^{(1/2)}_{m,k}=q^{k/2}(-q;q)_\infty J_k(w\one).
\]
The factor $q^{-k/2}$ in \eqref{eq:confhalf} cancels the displayed power of $q$. \Cref{prop:rhs-modular} proves modularity, and \eqref{eq:halfshift} supplies the negative half-shift sector.
\end{proof}

\section{Remaining classification and induction problems}\label{sec:open}

Three questions remain beyond the proved range.
\begin{enumerate}
\item The necessity direction of \cref{conj:odd} is open from rank eleven onward. A uniform argument must force the adjacent-pair relations and the ordered Boolean chain without a separate elimination in each rank.
\item The companion indices $7\le k\le m-1$ are unproved, beginning with $k=7$ in rank seventeen. The boundary inverse is known at every index (\cref{lem:inverse,thm:folding}); the missing step is the reduction on the Nahm side.
\item In even rank, sufficiency is open in rank six for five of the nine pairs of \cref{thm:r6}, and both directions are open from rank eight onward. \Cref{conj:even} below gives the expected answer, and \cref{conj:evenwron} reduces sufficiency to explicit theta-determinant identities, as \cref{conj:wron} does in odd rank. A related question is whether a number of cusp conditions independent of the rank suffices; ranks two, three, five, six, seven and nine need three coefficients, while rank four needs four.
\end{enumerate}

\subsection{Even ranks}
Ranks two, four and six suggest the following description.

\begin{conjecture}[Even-rank classification]\label{conj:even}
Let $r=2m$ with $m\ge1$. Then $\M_r$ contains the following $2m+3$ pairs, and for $m\ne2$ it consists exactly of them.
\begin{enumerate}
\item For integers $0\le x_0\le x_1\le\cdots\le x_{m-1}\le1$ with $k=\sum_ix_i$,
\[
B=(x_0,\,-x_1,\,x_1,\,\dots,\,-x_{m-1},\,x_{m-1},\,0),\qquad C=-\frac{r(2r+1)}{48(r+2)}+\frac{k(r+1-k)}{2(r+2)}.
\]
\item For integers $1\ge y_1\ge\cdots\ge y_{m-1}\ge0$ with $j=\sum_iy_i$,
\[
B=(y_1,\,-y_1,\,\dots,\,y_{m-1},\,-y_{m-1},\,0,\,\tfrac12),\qquad C=-\frac{r(r-4)}{24(r+2)}+\frac{j(r-j)}{2(r+2)}.
\]
\item $B=(1,-1,\dots,1,-1,1,-\tfrac12)$, with $m-1$ pairs $(1,-1)$, and $C=r/12$.
\item The vector of (iii) with $(1,-1)$ added to its first two coordinates, and $C=C_{m-1}+1$, where $C_{m-1}$ is the scalar in (ii) with $j=m-1$.
\end{enumerate}
For $m=2$, \cref{thm:r4} shows that $\M_4$ consists of these seven pairs together with the exceptional pair $\bigl((\frac32,-\frac32,\frac12,0),\frac9{16}\bigr)$.
\end{conjecture}

The sector (i) has the same scalar as \cref{conj:odd} with $z=0$. \Cref{thm:r2,thm:r4,thm:r6} agree with the conjecture in ranks two, four and six; in rank six only necessity and two of the nine sufficiency statements are proved. In ranks two and six the pair (iv) satisfies $F_{r,B_{\mathrm{(iv)}},0}=2q^{-1}F_{r,B_{\mathrm{(ii)},j=m-1},0}$, exactly in rank two and to order $q^{24}$ in rank six, which explains the shift of $C$ by one.

\subsection{An even-rank Wronskian formula}
The sums in \cref{conj:even} appear to satisfy theta-determinant identities parallel to \eqref{eq:wron}, with an alternating character on the lattice and without the factor $\prod_iv_i$. Let $r=2m$, $L=2m+2$ and $\nu=m(2m-1)$, and put
\[
\chi_0(v)=\prod_{1\le i<j\le m}(v_i^2-v_j^2),\qquad \epsilon_\tau(v)=(-1)^{\sum_i(v_i-\tau_i)/L}\quad(v\in\tau+L\Z^m).
\]
Define increasing tuples
\begin{align*}
\tau^{(0)}_k&=(\tfrac12,\tfrac32,\dots,m+\tfrac12)\setminus\{m+\tfrac12-k\}, & \tau^{(0)}_*&=(\tfrac12,\dots,m-\tfrac12),\\
\tau^{(1/2)}_j&=(0,1,\dots,m)\setminus\{m-j\}, & \tau^{(1/2)}_*&=(0,1,\dots,m-1),
\end{align*}
for $0\le k,j\le m$, and let
\[
R^{(z)}_\tau(q)=\frac{(-q^{1/2+z};q)_\infty^r}{(q;q)_\infty^\nu}\sum_{v\in\tau+L\Z^m}\epsilon_\tau(v)\,\frac{\chi_0(v)}{\chi_0(\tau^{(z)}_*)}\,q^{\frac{\|v\|^2-\|\tau\|^2}{2L}}.
\]

\begin{conjecture}[Even-rank Wronskian identity]\label{conj:evenwron}
Let $B^{\mathrm{(i)}}_k$ be the vector of \cref{conj:even}(i) with $\sum_ix_i=k$, let $B^{\mathrm{(ii)}}_j$ ($0\le j<m$) be the vector of \cref{conj:even}(ii) with $\sum_iy_i=j$, and let $B^{\mathrm{(ii)}}_m$ be the vector of \cref{conj:even}(iii). Then
\[
F_{2m,B^{\mathrm{(i)}}_k,0}=R^{(0)}_{\tau^{(0)}_k},\qquad F_{2m,B^{\mathrm{(ii)}}_j,0}=c_j\,R^{(1/2)}_{\tau^{(1/2)}_j},\qquad c_j=\begin{cases}1,&j<m,\\2,&j=m.\end{cases}
\]
Moreover, the vector $B^{\mathrm{(iv)}}$ of \cref{conj:even}(iv) satisfies $F_{2m,B^{\mathrm{(iv)}},0}=2q^{-1}F_{2m,B^{\mathrm{(ii)}}_{m-1},0}$.
\end{conjecture}

The cases $B^{\mathrm{(i)}}_0=0$ and $B^{\mathrm{(ii)}}_0=(0,\dots,0,\frac12)$ are the principal and last-coordinate half-shift formulas of \cite{CTW26} for even rank; the alternating character is their sign $\varepsilon=(-1)^{r+1}$. In rank two the last relation is proved in \cref{thm:r2}. In odd rank the analogous relation \eqref{eq:halfshift} follows from the simultaneous shift of \cref{lem:simshift}, which comes from the odd-rank reduction of Shi and Wang. Their even-rank reduction \cite[(1.10)]{SW26}, which holds with $x_0=x_1^{-1}$, admits no shift of this kind: the factors $x_0^{-n_1}$ and $(x_2/x_0)^{n_2}$ leave an uncancelled term $q^{hn_2}$. A proof of the last relation, and of \cref{conj:evenwron} beyond the principal cases, therefore seems to require the Hall--Littlewood identities of \cite{BW15} at non-principal specializations, as in \cref{sec:theta}.

\begin{proposition}\label{prop:evenwron}
For each tuple $\tau$ above, $q^{C}R^{(z)}_\tau$ is a modular function on a congruence subgroup, where $C$ is the scalar attached to the corresponding pair in \cref{conj:even}. Consequently, \cref{conj:evenwron} implies that every pair of \cref{conj:even} is modular, and in particular that $\M_6$ consists of exactly the nine pairs of \cref{thm:r6}.
\end{proposition}

\begin{proof}
For rational $a$, the series $\theta^-_{a,L}(q)=\sum_{v\in a+L\Z}(-1)^{(v-a)/L}q^{v^2/(2L)}$ has weight $1/2$ on a congruence subgroup. As in the proof of \cref{prop:rhs-modular}, expanding the Wronskian $\det(\mathcal D^{i-1}\theta^-_{\tau_j,L})_{i,j=1}^m$ and using the Vandermonde determinant in the variables $v_j^2$ gives a constant multiple of $\sum_{v\in\tau+L\Z^m}\epsilon_\tau(v)\chi_0(v)q^{\|v\|^2/(2L)}$, which has weight $m/2+m(m-1)=\nu/2$. The eta-quotient factors are those of \cref{prop:rhs-modular}, and the weights cancel. The power of $q$ that normalizes the product is
\[
\frac{\|\tau\|^2}{2L}-\frac\nu{24}-\begin{cases}\frac r{48},&z=0,\\[2pt]-\frac r{24},&z=\frac12,\end{cases}
\]
and direct evaluation of $\|\tau\|^2$ shows that it equals the scalar of \cref{conj:even} in sectors (i)--(iii). The relation for $B^{\mathrm{(iv)}}$ multiplies by $2q^{-1}$, which shifts the scalar by one, as in (iv). Together with \cref{thm:r6}, this gives the last assertion.
\end{proof}

\subsection{Even rank: the first companions}\label{subsec:evenproof}
The argument of \cref{sec:theta,sec:boundary,sec:cert} transfers to the sector (i) of \cref{conj:even}, with the type-C identity of Bartlett and Warnaar replaced by its type-B companion. Throughout, $r=2m$, $L=2m+2$, $\nu=m(2m-1)$ and $w=q^{1/2}$. For generic $\hat x=(x_1,\dots,x_{m-1})$ put $x=(\hat x,1)$; thus $x_m=1$ is frozen. Define
\[
\HL^B_m(t;\hat x;q)=\sum_{\lambda:\ \lambda_1\le2}q^{|\lambda|/2}\,t^{m_1(\lambda)}\,P'_\lambda(x_1,x_1^{-1},\dots,x_{m-1},x_{m-1}^{-1},1;q),
\]
a sum over Hall--Littlewood functions in $2m-1$ letters, in which the letter $1$ occurs once. Let
\begin{gather*}
\mu^-_x[f]=\sum_{\ell\in\Z}(-1)^\ell q^{L\ell^2/2}x^{L\ell}f(xq^\ell),\\
\Phi^+_a(x)=\mu^-_x[y^a+y^{-a}],\qquad \Psi^+_a(x)=\mu^-_x\Bigl[\frac{y^a+y^{-a}}{1+w^2+w(y+y^{-1})}\Bigr],
\end{gather*}
let $\delta^+_\varepsilon(x)=\det\bigl(x_i^{j-1+\varepsilon}+x_i^{-(j-1+\varepsilon)}\bigr)_{i,j=1}^m$ for $\varepsilon\in\{0,\frac12\}$, and
\[
P^B_t(\hat x)=\frac{(-wt;q)_\infty\prod_{i<m}(-wtx_i,-wt/x_i;q)_\infty}{(q;q)_\infty^m\prod_{i<m}(qx_i^{\pm2},qx_i^{\pm1};q)_\infty\prod_{1\le i<j<m}(qx_i^{\pm}x_j^{\pm};q)_\infty},
\]
where $(qx_i^{\pm2};q)_\infty=(qx_i^2,qx_i^{-2};q)_\infty$, and $(qx_i^\pm x_j^\pm;q)_\infty$ denotes the product over the four choices of signs. The frozen letter contributes the single factor $(-wt;q)_\infty$ and the cross factors $(qx_i^{\pm1};q)_\infty$.

\begin{proposition}[Even boundary formulas]\label{prop:evenbd}
\begin{enumerate}
\item For rational $c$,
\[
F_{2m,(0,\dots,0,c),0}(q)=(-q^{c+1/2};q)_\infty\,\HL^B_m(q^c;\one;q).
\]
\item For generic $\hat x$,
\begin{align*}
\HL^B_m(1;\hat x;q)&=P^B_1(\hat x)\,\frac{\det(\Phi^+_{j-1/2}(x_i))_{i,j=1}^m}{\delta^+_{1/2}(x)},\\
\HL^B_m(q;\hat x;q)&=(1+w)\,P^B_1(\hat x)\,\frac{\det(\Psi^+_{j-1/2}(x_i))_{i,j=1}^m}{\delta^+_{1/2}(x)},\\
\HL^B_m(w;\hat x;q)&=P^B_w(\hat x)\,\frac{\det(\Phi^+_{j-1}(x_i))_{i,j=1}^m}{\delta^+_{0}(x)}.
\end{align*}
\end{enumerate}
\end{proposition}

\begin{proof}[Proof sketch]
Part (i) is the case $n=2m-1$, $u=1$ of \cite[Theorem~3.7]{BW15}, after Euler's identity on the last coordinate as in \cref{lem:boundary} (equivalently \cite[Lemma~2.1]{CTW26}). For (ii), take level one, let $b\to\infty$ and put $c=-w/t$ in \cite[(5.6b)]{BW15}, which has $x=(\hat x,1)$, $K=n+\frac12$ and the factor $\Delta_B(-xq^r)/\Delta_B(-x)$. The one-variable factor becomes
\[
(-1)^{r}\,t^{r}\,q^{(m+1)r^2+(\frac12-m)r}\,x^{Lr}\,\frac{(-wx/t;q)_r}{(-wtx;q)_r}.
\]
Write $\Delta_B(y)=\prod_i(1-y_i)\prod_{i<j}(y_i-y_j)(y_iy_j-1)$ as $(-1)^m\prod_iy_i^{m-1/2}$ times the type-B determinant $\det\bigl(y_i^{m-j+1/2}-y_i^{-(m-j+1/2)}\bigr)$. At $y=-xq^r$ the monomial $\prod_i(-x_iq^{r_i})^{m-1/2}$ absorbs the factor $q^{(m-1/2)r}$, and for half-integral $a$ one has $(-y)^a-(-y)^{-a}=i^{2a}(y^a+y^{-a})$. Multilinearity in the rows turns the sum over $r\in\Z^m$ into the cosine determinant with entries $\Phi^+$; the constants cancel against the same computation at $r=0$, which produces $\delta^+_{1/2}$. This gives the case $t=1$. At $t=q$, the two telescoping factors combine with $P^B_q/P^B_1$ as in the proof of \cref{lem:boundary}; for the frozen letter the combination leaves the constant $1+w$. At $t=w$, the quotient $(1+x)/(1+xq^\ell)$ cancels a factor of the type-B denominator row by row, and the remaining determinant has integral indices.
\end{proof}

\begin{lemma}[Type-B tridiagonal inverse]\label{lem:evenB}
For $0\le k\le m$ and $\tau=\tau^{(0)}_k$ put
\[
J^B_k(\hat x)=P^B_1(\hat x)\,\frac{\det(\Phi^+_{\tau_j}(x_i))_{i,j=1}^m}{\delta^+_{1/2}(x)}.
\]
Then
\[
\sum_{j=0}^kA_{k,j}(q)\,\HL^B_m(q^j;\hat x;q)=q^{k^2/2}J^B_k(\hat x)\qquad(0\le k\le m).
\]
In particular $\HL^B_m(1;\hat x;q)-\HL^B_m(q;\hat x;q)=wJ^B_1(\hat x)$.
\end{lemma}

\begin{proof}
Put $d=m+1$ and index the functions by $a\in\{\frac12,\frac32,\dots,d-\frac12\}$. As in \cref{lem:bdexp}, $\Phi^+_a=(1+t^2)\Psi^{+[1]}_a+t(\Psi^{+[1]}_{a-1}+\Psi^{+[1]}_{a+1})$ with $t=w$, and similarly at every level $t=wq^{j-1}$. Two boundary effects change. At the bottom, the cosine symmetry gives $\Psi^+_{-1/2}=\Psi^+_{1/2}$, so the first diagonal entry of the tridiagonal matrix becomes $1+t+t^2$ and its determinant becomes $1+t$. At the top, the reindexing $\mu^-_x[f(y)]=-q^{L/2}\mu^-_x[y^Lf(qy)]$ applied to $f(y)=y^{-A}/(1+t/y)$ with $A=d+\frac12=(L+1)/2$ gives $\Psi^+_{d+1/2}+t\Psi^+_{d-1/2}=0$, so the last diagonal entry is $1$ as before. For $t=w$, the complementary-minor formula and Cauchy--Binet give
\[
\det\bigl(\Psi^+_{1/2},\dots,\Psi^+_{m-1/2}\bigr)=\frac{\det(\Phi^+_{1/2},\dots,\Phi^+_{m-1/2})-w\det(\Phi^+_{1/2},\dots,\Phi^+_{m-3/2},\Phi^+_{m+1/2})}{1+w},
\]
and \cref{prop:evenbd} gives the case $k=1$. For general $k$, the products $V_j$ of the modified matrices have determinant $(-w;q)_j$, which cancels in the complementary-minor formula, so \eqref{eq:bdexp} holds with the same coefficients $(V_j)_{d-h,d}$. For $j\le m$ no matrix path from the last coordinate reaches the modified first diagonal entry, so these coefficients agree with \cref{prop:coeffs}, and \cref{lem:inverse} applies verbatim.
\end{proof}

At $\hat x=\one$ the confluent limit is computed as in \cref{lem:confluent}, with the Vandermonde determinant in the variables $\log^2x_i$: since $\Phi^+_a(x)=\sum_{v\in a+L\Z}\epsilon_a(v)q^{(v^2-a^2)/(2L)}(x^v+x^{-v})$,
\begin{equation}\label{eq:evenconf}
J^B_k(\one)=\frac{(-w;q)_\infty^{2m-1}}{(q;q)_\infty^\nu}\sum_{v\in\tau+L\Z^m}\epsilon_\tau(v)\frac{\chi_0(v)}{\chi_0(\tau^{(0)}_*)}q^{\frac{\|v\|^2-\|\tau\|^2}{2L}},\qquad\tau=\tau^{(0)}_k.
\end{equation}

\begin{theorem}[Even companions in sector (i)]\label{thm:evencomp}
Let $m\ge2$ and $1\le k\le\min(6,m-1)$. Then the identity of \cref{conj:evenwron} holds for $B^{\mathrm{(i)}}_k$, that is, $F_{2m,B^{\mathrm{(i)}}_k,0}=R^{(0)}_{\tau^{(0)}_k}$. Consequently $\bigl(B^{\mathrm{(i)}}_k,C\bigr)\in\M_{2m}$ with the scalar $C$ of \cref{conj:even}(i).
\end{theorem}

\begin{proof}
The vector $B^{\mathrm{(i)}}_k$ is zero except for $k$ pairs $(-1,1)$ in positions $(2m-2i,2m-2i+1)$, $1\le i\le k$. For $k=1$, the contiguity computation leading to \eqref{eq:contig2} applies verbatim, since $T_{2m}(e_{2m-1}+e_{2m})=e_{2m-1}-e_{2m-2}$; with $c=0$ and \cref{prop:evenbd}(i) it gives $F_{2m,B^{\mathrm{(i)}}_1,0}=q^{-1/2}(-w;q)_\infty(\HL^B_m(1;\one;q)-\HL^B_m(q;\one;q))$. For $2\le k\le6$, the reductions of \cref{lem:k2red,lem:certred} at $v=0$, that is $u=1$, involve only the last $2k+1$ coordinates, and their proofs use only the local form of the tadpole quadratic form near the last coordinate. They therefore hold in rank $2m$ whenever $2m\ge2k+1$, and give
\[
q^{k^2/2}F_{2m,B^{\mathrm{(i)}}_k,0}=\sum_{j=0}^kA_{k,j}(q)\,(-w;q)_j\,F_{2m,je_{2m},0}=(-w;q)_\infty\sum_{j=0}^kA_{k,j}(q)\,\HL^B_m(q^j;\one;q).
\]
(For $k=5,6$ this uses the certificates at $u=1$.) \Cref{lem:evenB} and \eqref{eq:evenconf} identify the right-hand side with $R^{(0)}_{\tau^{(0)}_k}$, and \cref{prop:evenwron} gives modularity.
\end{proof}

\begin{corollary}\label{cor:r6more}
The second and third pairs of \cref{thm:r6} are modular. Thus four of the nine pairs in rank six are proved modular, and for every $m\ge7$ the set $\M_{2m}$ contains at least the eight pairs of \cref{conj:even}(i) with $k\le6$ and (ii) with $j=0$.
\end{corollary}

The remaining even sectors, namely (i) with $k=m$, (ii) with $j\ge1$, (iii) and (iv), contain the shift $v=\frac12$ of the boundary family $B=(2v,-2v,\dots,2v,c)$. For $v\ne0$, \cite[Theorem~3.7]{BW15} gives Hall--Littlewood functions in the letters $(u,u^{-1},\dots,u^{-1},u)$ with $u=q^v$, in which the unpaired letter is $u$ rather than $1$. This is not of the form covered by \cite[(5.6b)]{BW15}, and a determinant formula for it is the missing ingredient for these sectors. In plethystic terms the obstruction is the unpaired letter: for $u=w$ one has $p_k[X/(1-q)]=\bigl((m-1)(w^k+w^{-k})+w^k\bigr)/(1-q^k)$, which is not of the form produced by the alphabets $(x^\pm)$ or $(\hat x^\pm,1)$.

\begin{remark}\label{rem:evenlead}
The constant terms of the right-hand sides in \cref{conj:evenwron} are
\[
\frac{\chi_0(\tau^{(0)}_k)}{\chi_0(\tau^{(0)}_*)}=\binom{2m-1}{k}-\binom{2m-1}{k-2},\qquad c_j\,\frac{\chi_0(\tau^{(1/2)}_j)}{\chi_0(\tau^{(1/2)}_*)}=\binom{2m}{j},
\]
by cancellation in the products defining $\chi_0$. These are the even analogues of \cref{prop:leading}. The factor $c_m=2$ is exactly what is needed for $\binom{2m}{m}$. The Nahm sums have the same constant terms in every case computed (ranks two to fourteen), and a path-counting proof as in \cref{prop:leading} should be possible.
\end{remark}


\begin{remark}\label{rem:evenwron-evidence}
The identities of \cref{conj:evenwron} not covered by \cref{thm:evencomp} or \cite{CTW26} agree with the Nahm sums as $q$-series, in exact integer arithmetic, through $q^{40}$ in rank six, through $q^{15}$ in ranks eight, ten and twelve, and through $q^{10}$ in rank fourteen. The relation for $B^{\mathrm{(iv)}}$ agrees through $q^{20}$ in ranks four to eight. In rank four the corresponding sums are known to be modular by \cite{SW26,SW26b}. This evidence is not used in any proof.
\end{remark}

\begin{remark}\label{rem:highrank}
The conjectures have been tested against the cusp conditions beyond the classified ranks. For each of the $11$ and $13$ pairs predicted by \cref{conj:even} in ranks eight and ten, and for each of the $18$ and $21$ pairs predicted by \cref{conj:odd} in ranks eleven and thirteen, the procedure of \cref{sec:cusp} gives $\kappa_2=\kappa_3=0$ and $c_1-\gamma_0=C$ to within $10^{-10}$ in floating-point arithmetic. Random half-integral shifts give values between $10^{-2}$ and $10^{2}$. In rank eleven, replacing the ordered chain $x_1\le\cdots\le x_5$ by unordered patterns such as $(1,0,0,0,0)$ or $(1,0,1,0,1)$ gives $|\kappa_2|>2\cdot10^{-2}$; thus the ordering in \cref{conj:odd} is not an artefact of the ranks used to formulate it. These computations are evidence only.
\end{remark}

\begin{remark}\label{rem:numerics}
As a check on normalizations only, identity \eqref{eq:wron} has been compared coefficientwise, in exact integer arithmetic, through $q^{15}$ for $(m,k)=(8,7)$ and $(9,8)$ in both sectors and for $(m,k)=(9,7)$ and $(10,9)$ in the sector $z=0$. These cases lie outside the proved range. Such comparisons are not used anywhere as proofs.
\end{remark}

\subsection{The missing boundary reduction}
Set $\HL_{m,j}=\HL_m(q^j;\one;q)$. The reduction in the sector $z=0$ is
\begin{equation}\label{eq:missing}
\frac{S^{(0)}_{m,k}}{(-w;q)_\infty}=q^{-k^2/2}\sum_{j=0}^kA_{k,j}(q)\,\HL_{m,j}.
\end{equation}
\Cref{thm:basic,thm:k2,thm:k3to6}, together with \cref{lem:inverse}, prove this identity for $0\le k\le\min(6,m)$ and for $k=m$. The remaining interior indices are conjectural. The local formulation below would prove the more general reduction \eqref{eq:certred} by induction.

\subsection{A local induction problem for the remaining companions}
The boundary relation at index $m+1$ supplies the initial value of a more demanding contiguity identity, which we now state explicitly; the basis theorem alone does not prove it.

Retain $u=q^v$, $w=q^{1/2}$ and the polynomials $R_s(e)$ defined above, now for every $s\ge0$. For $k\ge1$, put $r_k=2k-1$, and let $b^{(k)}_v\in\Q^{r_k}$ consist of $k-1$ pairs $(-2v,2v)$ followed by $-v$. For integers $N,s\ge0$, define
\begin{equation}\label{eq:moment}
\mathfrak M^{(s)}_{k,N}(v;q)=\sum_{j=0}^sA_{s,j}(q)\,(-w/u;q)_j\,F_{r_k,\,b^{(k)}_v-Ne_1+je_{r_k},\,0}(q).
\end{equation}
For $k\ge2$, \cref{cor:bdrel} proves $\mathfrak M^{(k)}_{k,0}=0$, and Euler's identity proves the same for $k=1$. The proposed extension is
\begin{equation}\label{eq:extension}
\mathfrak M^{(k)}_{k,N}(v;q)=\frac{q^{k-1/2}}u\,(q^{-N}-1)\,\mathfrak M^{(k-1)}_{k,N-1}(v;q)\qquad(N\ge1).
\end{equation}
Both moments here have rank $2k-1$. We do not establish this identity for arbitrary $k$.

The following proposition gives its precise role. For $n=(n_1,\dots,n_{2k})\in\Z_{\ge0}^{2k}$ and indeterminates $a=(a_1,\dots,a_{2k})$, set
\[
W_k(n;u)=\frac{q^{\sum_{i=1}^{2k-1}n_i^2+n_{2k}^2/2-\sum_{i=1}^{2k-1}n_in_{i+1}}\,u^{2\sum_{i=1}^{2k-1}(-1)^{i+1}n_i-n_{2k}}}{\prod_{i=1}^{2k}(q;q)_{n_i}},\qquad T_k(a)=\frac{a_1a_3\cdots a_{2k-1}}{a_2a_4\cdots a_{2k-2}},
\]
where an empty product is $1$. For a Laurent polynomial $P$ in $2k$ variables, write
\[
\mathcal Z_k(X;P)=\sum_{n\ge0}W_k(n;u)\,X^{n_1}P(q^{n_1},\dots,q^{n_{2k}}).
\]
Positive definiteness of $T_{2k}$ makes this an entire function of $X$. The local boundary identity is
\begin{equation}\label{eq:local}
\mathcal Z_k(X;R_k(a_{2k}))=q^{k^2/2}u^kX\,\mathcal Z_k(X;T_k).
\end{equation}

\begin{proposition}[Induction criterion]\label{prop:induction}
Identity \eqref{eq:local} holds for $k=1$ and every rational $v$. For $k\ge2$, assume that it holds at $k-1$ for every $X$. Then \eqref{eq:local} at $k$ is equivalent to \eqref{eq:extension} for every $N\ge1$. If these identities hold through index $k$, they imply \eqref{eq:certred} at that index for every $m\ge k$.
\end{proposition}

\begin{proof}
The coefficient of $X^N$ on the left of \eqref{eq:local} is
\[
\frac{q^{N^2}u^{2N}}{(q;q)_N}\,\mathfrak M^{(k)}_{k,N}.
\]
For $N\ge1$, fix $n_1=N-1$ on the right. The remaining insertion is the full companion of index $k-1$ on the coordinates $n_2,\dots,n_{2k}$, with an additional first-coordinate shift $-(N-1)e_1$. Apply the local identity at $k-1$, first fixing $n_2$ and then summing over it. The resulting coefficient is
\[
\frac{q^{N^2-N+k-1/2}u^{2N-1}}{(q;q)_{N-1}}\,\mathfrak M^{(k-1)}_{k,N-1}.
\]
Equality of these coefficients is precisely \eqref{eq:extension}. The constant coefficients vanish on both sides, since $\mathfrak M^{(k)}_{k,0}=0$. Since both sides are entire in $X$, this proves the equivalence.

For the base case, put $c=w/u$. Euler's identity gives
\[
\mathfrak M^{(0)}_{1,N}=(-cq^{-N};q)_\infty,\qquad \mathfrak M^{(1)}_{1,N}=c(q^{-N}-1)(-cq^{1-N};q)_\infty,
\]
and substitution into the coefficient calculation verifies \eqref{eq:local} for $k=1$, including its constant term.

Finally, fix the coordinate $n_{r-2k}$ of a Nahm sum of rank $r=2m+1$ and set $X=q^{-n_{r-2k}}$. The factor $XT_k$ is exactly the insertion for its last $k$ companion pairs. Summing \eqref{eq:local} over the fixed coordinate and all preceding coordinates, which is justified by absolute convergence, yields
\[
q^{kv+k^2/2}G_{m,k}(v;q)=\sum_{j=0}^kA_{k,j}(q)\,(-w/u;q)_j\,F_{r,B_v+je_r,0}(q).
\]
\Cref{lem:boundary} then gives \eqref{eq:certred}.
\end{proof}

The telescoping certificates of \cref{sec:cert} prove the local identity through $k=4$ for arbitrary rational $v$, and through $k=6$ for $v=0,\frac12$; consequently they also prove \eqref{eq:extension} in those ranges. For larger indices, the value at $N=0$ alone does not determine the values at $N\ge1$, and establishing that extension is the remaining step in the proposed induction. \Cref{prop:recurrence} is an identity among the polynomials $R_k$ and does not supply this extension.


\appendix
\section{Exact obstruction data}\label{app:data}

These data implement the cusp conditions of \cref{sec:cusp}. The larger polynomial expressions are recorded in the accompanying data (\cref{app:verification}).

\subsection{Rank three}\label{app:r3}
For $A=T_3$, write $s=\sqrt5$. Direct substitution gives
\[
Q=\Bigl(\frac{3-s}2,\ s-2,\ \frac{3-s}4\Bigr),\qquad \xi=\Bigl(\frac{s-1}2,\ \frac{s-1}4,\ s-2\Bigr),
\]
and
\[
M=\begin{pmatrix}-\frac32+\frac{9s}{10}&-1+\frac{3s}5&\frac12-\frac s{10}\\[2pt] -1+\frac{3s}5&\frac{2s}5&\frac12+\frac s{10}\\[2pt] \frac12-\frac s{10}&\frac12+\frac s{10}&\frac12+\frac{2s}5\end{pmatrix}.
\]
These saddle-point values and the first-order condition already occur in \cite[Section~4.1]{MW24}. The recurrence of \cref{sec:cusp} is used with $N=3$, so the coefficients lie in $\Q(\sqrt5)[b_1,b_2,b_3]$. The first-order calculation yields
\begin{equation}\label{eq:r3first}
C=C_0(B)+\frac{\sqrt5}{20}P(B),
\end{equation}
where
\begin{align*}
80C_0(B)&=-60b_1^2-80b_1b_2+40b_1b_3-72b_1+40b_2b_3-68b_2+20b_3^2+8b_3-7,\\
P(B)&=9b_1^2+12b_1b_2-2b_1b_3+8b_1+4b_2^2+2b_2b_3+7b_2+4b_3^2-2b_3.
\end{align*}
Since $B$ and $C$ are rational, $P(B)=0$ and $C=C_0(B)$. For rational $B$, both the rational and the $\sqrt5$ components of $\kappa_2$ and $\kappa_3$ must also vanish; these are four further polynomial equations.

After applying \eqref{eq:r3coords} and clearing denominators, the five equations are $E_0=\cdots=E_4=0$, where
\begin{align*}
E_0&=9u^2-6uv+10uz+8u+v^2-v,\\
E_1&=7200u^3-6000u^2v+8400u^2z+16728u^2+1200uv^2-3600uvz-5832uv\\
&\qquad+2400uz^2+8640uz+9336u-48v^2+48v-100z^3+25z,\\
E_2&=-812u^3+672u^2v-960u^2z-1866u^2-144uv^2+420uvz+654uv\\
&\qquad-300uz^2-960uz-1039u-4v^3+6v^2-2v.
\end{align*}
The two quartics $E_3,E_4$, and all certificate multipliers, are recorded as exact rational expressions in the accompanying data. The precise normalization is
\[
E_0=P(B),\qquad \kappa_2=\frac{E_1}{2400}+\frac{\sqrt5\,E_2}{600},\qquad \kappa_3=\frac{E_3}{2000}+\frac{\sqrt5\,E_4}{6000},
\]
where the $E_i$ are evaluated using \eqref{eq:r3coords}.

\begin{lemma}[Exact elimination]\label{lem:r3elim}
In $\Q[u,v,z]$, the ideal generated by $E_0,\dots,E_4$ equals $\bigl(u,\ v(v-1),\ z(4z^2-1)\bigr)$.
\end{lemma}

\begin{proof}
The certificate gives rational polynomials $A_0,\dots,A_4$ satisfying
\begin{equation}\label{eq:r3cert}
\sum_{i=0}^4A_iE_i=u,
\end{equation}
each product having total degree at most six. Expanding the left-hand side with exact rational coefficients proves \eqref{eq:r3cert}. Consequently $u$ belongs to the obstruction ideal. Next,
\[
E_0\big|_{u=0}=v(v-1),\qquad E_1\big|_{u=0}=-48v(v-1)-25z(4z^2-1),
\]
so the other two displayed generators belong to the ideal as well. Conversely, each $E_i$ reduces to zero under $u=0$, $v^2=v$ and $z^3=z/4$, as verified by direct polynomial division. This proves the reverse containment.
\end{proof}

\subsection{Rank four}\label{app:r4}
Write $s=\sqrt3$. The unique saddle in $(0,1)^4$ is
\[
Q=\Bigl(\frac{s-1}2,\ \frac{3-s}6,\ \frac{2s-3}3,\ \frac{2-s}2\Bigr),
\]
as direct substitution into the four saddle equations confirms. The recurrence of \cref{sec:cusp}, with $r=N=4$, computes $c_1,c_2,c_3,c_4$ over $\Q(\sqrt3)$. The additional logarithmic coefficient is
\[
\kappa_4=c_4-c_1c_3-\tfrac12c_2^2+c_1^2c_2-\tfrac14c_1^4,
\]
and modularity requires $\kappa_4=0$ as well as \eqref{eq:cuspcond}. The first-order condition takes the form $C=C_4(B)+\sqrt3R_4(B)$, where $C_4,R_4\in\Q[b_1,b_2,b_3,b_4]$ and
\begin{align*}
72C_4(B)&=24b_1^2+18b_1b_2+12b_1b_3+6b_1b_4-6b_1+20b_2^2+28b_2b_3+18b_2b_4-4b_2\\
&\qquad+32b_3^2+24b_3b_4-4b_3+24b_4^2+6b_4-9,\\
72R_4(B)&=-3b_1^2+4b_1b_2+4b_1b_3+6b_1b_4+3b_1+7b_2^2+12b_2b_3+14b_2b_4-b_2\\
&\qquad+8b_3^2+32b_3b_4-4b_3+24b_4^2-12b_4.
\end{align*}
Since $B$ and $C$ are rational, $R_4(B)=0$ and $C=C_4(B)$. The rational and $\sqrt3$ components of each of $\kappa_2,\kappa_3,\kappa_4$ also vanish. Let $E^{(4)}_0,\dots,E^{(4)}_6$ be the primitive integer polynomials obtained from these seven rational polynomials in the order $R_4$, then the rational and irrational components for the indices $2,3,4$. Here \emph{primitive} means that denominators are cleared and the common integer factor is removed. Their degrees are $2$, $3$, $3$, $4$, $4$, $5$, $5$, and their exact expressions are recorded in the accompanying data.

Let $J_4=(E^{(4)}_0,\dots,E^{(4)}_6)\subset\Q[b_1,b_2,b_3,b_4]$, and define
\begin{align*}
G_1&=b_1^2-b_1-2b_2b_4+b_2-b_3^2+2b_3+2b_4^2-b_4,\\
G_2&=b_1b_2+b_1+2b_2b_4-b_2+b_3^2-2b_3-2b_4^2+b_4,\\
G_3&=2b_1b_3-2b_1-2b_2b_4-b_2-2b_3^2+b_3+2b_4^2-b_4,\\
G_4&=b_4(b_1+b_2),\\
G_5&=b_2^2-2b_2b_4+2b_2-b_3^2+2b_3+2b_4^2-b_4,\\
G_6&=2b_2b_3+2b_2b_4-b_2+2b_3^2-b_3-2b_4^2+b_4,\\
G_7&=4b_2b_4^2-b_2+4b_3^2-5b_3+2b_4^2-b_4,\\
G_8&=b_3(2b_3-1)(b_3-1),\\
G_9&=b_4(2b_3+2b_4-1),\\
G_{10}&=b_4(2b_4-1)(2b_4+1).
\end{align*}

\begin{lemma}[Rank-four elimination]\label{lem:r4elim}
The polynomials $G_1,\dots,G_{10}$ form a Gr\"obner basis of $J_4$ for the lexicographic order $b_1>b_2>b_3>b_4$. The ideal is radical, and its complex zero set consists of the eight rational vectors of \cref{thm:r4}.
\end{lemma}

\begin{proof}
The accompanying certificate gives rational polynomials $A_{ji}$ such that
\[
G_j=\sum_{i=0}^6A_{ji}E^{(4)}_i\qquad(1\le j\le10),
\]
each product $A_{ji}E^{(4)}_i$ having total degree at most nine. These ten identities hold by exact polynomial multiplication. Conversely, division by the $G_j$ reduces each $E^{(4)}_i$ to zero. All $\binom{10}2=45$ Buchberger $S$-polynomials reduce to zero, which proves the Gr\"obner-basis assertion.

The monomials not divisible by the leading monomials of this basis are
\[
1,\ b_4,\ b_4^2,\ b_3,\ b_3^2,\ b_2,\ b_2b_4,\ b_1.
\]
They form a basis of the quotient ring, which therefore has dimension eight over $\Q$. Direct substitution gives eight distinct rational zeros, namely the vectors of the theorem. Hence there are no further complex zeros and no nontrivial multiplicities, and $J_4$ is radical.
\end{proof}

\subsection{Rank five}\label{app:r5}
Let $t$ be the unique root in $(9/100,1/10)$ of
\[
8t^3-20t^2+12t-1=0;
\]
equivalently, $t=1-\cos(\pi/7)$. The polynomial is irreducible over $\Q$, and direct substitution gives the saddle
\[
Q=\bigl(-4t^2+4t,\ 2t,\ 12t^2-20t+2,\ -8t^2+12t-1,\ t\bigr),
\]
each coordinate of which lies in $(0,1)$. The first three expansion coefficients are computed over $K=\Q(t)$ by the recurrence of \cref{sec:cusp}. Write
\[
c_1-\gamma_0=C_0+tC_1+t^2C_2,\qquad \kappa_j=K_{j,0}+tK_{j,1}+t^2K_{j,2}\quad(j=2,3),
\]
where all component polynomials have rational coefficients in $B$. For rational $B$ and $C$, modularity requires
\[
C_1=C_2=K_{2,0}=K_{2,1}=K_{2,2}=K_{3,0}=K_{3,1}=K_{3,2}=0,\qquad C=C_0.
\]
Clear denominators and remove common integer factors from these eight polynomials, in the displayed order. After the change of coordinates of \cref{thm:r5}, denote them by $E^{(5)}_0,\dots,E^{(5)}_7$. Their degrees are $2$, $2$, $3$, $3$, $3$, $4$, $4$, $4$. Their exact expressions are recorded in the accompanying data.

\begin{lemma}[Rank-five elimination]\label{lem:r5elim}
In $\Q[u,v,x,y,z]$, the obstruction ideal is
\[
\bigl(E^{(5)}_0,\dots,E^{(5)}_7\bigr)=\bigl(u,\ v,\ x^2-x,\ xy-x,\ y^2-y,\ 4z^3-z\bigr).
\]
It is radical and has exactly the nine rational zeros described in \eqref{eq:r5set}.
\end{lemma}

\begin{proof}
For each of the six generators $G_j$ on the right, the certificate records eight rational polynomials $A_{ji}$ such that
\[
G_j=\sum_{i=0}^7A_{ji}E^{(5)}_i,\qquad \deg\bigl(A_{ji}E^{(5)}_i\bigr)\le7.
\]
These identities hold by exact expansion. Conversely, every $E^{(5)}_i$ reduces to zero modulo the six displayed generators. These generators form a Gr\"obner basis for the lexicographic order $u>v>x>y>z$; all fifteen $S$-polynomials reduce to zero. The standard monomials are $\{1,x,y\}\cdot\{1,z,z^2\}$, so the quotient has dimension nine. The three choices of $(x,y)$ in \eqref{eq:r5set}, together with the three distinct roots of $4z^3-z$, give nine rational zeros. This proves radicality and excludes any further complex zeros.
\end{proof}

\section{Computations used in the proofs}\label{app:verification}

This appendix lists the computations on which the proofs depend. Statements proved by hand, in particular those of \cref{sec:theta,sec:boundary} and the evaluations in \cref{thm:r2}, use no computation. The data and the programs that carry out the computations below are available in the repository cited under Data availability.

\emph{Cusp polynomials.} In every rank, the obstruction polynomials are obtained from the procedure of \cref{sec:cusp} in exact arithmetic over $K=\Q(t)$, $t=1-\cos(\pi/(r+2))$. The saddle point is obtained from $Q_r=t$ by solving the saddle equations backwards, and these equations are then checked exactly in $K$. In rank nine, the polynomials are computed modulo $p$ at the five roots of $f$, as described in the proof of \cref{thm:r9}.

\emph{Ranks three, four and five.} The eliminations of \cref{lem:r3elim,lem:r4elim,lem:r5elim} are proved by rational ideal-membership certificates and Gr\"obner bases. Each identity is checked by exact rational expansion, and the reverse containments by exact polynomial division.

\emph{Ranks two, six, seven and nine.} These eliminations use \cref{lem:transfer}. For each rank, two computations are needed over $\F_p$: the quotient dimension of the affine obstruction ideal, and the finiteness of the quotient by the ideal of highest homogeneous parts. The primes are $p=1000003$ in ranks two, six and seven, and $p=1000033$ in rank nine. In rank nine, the local integrality at $p$ and the preservation of degrees under reduction are part of the proof of \cref{thm:r9}. In ranks two and six, the rational zeros are located on a finite grid and then verified by exact substitution; since their number equals the bound from \cref{lem:transfer}, no other zeros exist.

\emph{Companion certificates.} \Cref{lem:certred} relies on the certificate coefficients for $k=3,\dots,6$, and identity \eqref{eq:certid} is checked by exact expansion in the Laurent polynomial ring over $\Z$. The fifth and sixth companions use separate certificates at $u=1$ and $u=w$. For $k=6$ the certificate has been cleared of an explicitly nonzero polynomial denominator $D_{6,v}(w)$, and the identity theorem justifies cancellation across its zeros. \Cref{thm:evencomp} uses the certificates at $u=1$. The certificate for $k=2$ is printed in full in \eqref{eq:k2cert}.


\section*{Data availability}
The data and programs for the computations described in \cref{app:verification} are available at \url{https://github.com/ramanujamath/tadpole-nahm-modularity} (archived version: [DOI to be inserted]).

\section*{Acknowledgements}
[To be completed by the author(s).]

\begin{thebibliography}{99}

\bibitem{Andrews76}
G.~E.~Andrews, \emph{The Theory of Partitions}, Encyclopedia of Mathematics and its Applications~2, Addison-Wesley, 1976; reissued by Cambridge University Press, 1998.

\bibitem{BW15}
N.~Bartlett and S.~O.~Warnaar, \emph{Hall--Littlewood polynomials and characters of affine Lie algebras}, Adv. Math. \textbf{285} (2015), 1066--1105. \href{https://doi.org/10.1016/j.aim.2015.08.011}{doi:10.1016/j.aim.2015.08.011}; arXiv:1304.1602.

\bibitem{msolve}
J.~Berthomieu, C.~Eder and M.~Safey El Din, \emph{msolve: a library for solving polynomial systems}, in: Proceedings of the 2021 International Symposium on Symbolic and Algebraic Computation (ISSAC '21), ACM, 2021. \href{https://doi.org/10.1145/3452143.3465545}{doi:10.1145/3452143.3465545}.

\bibitem{CGZ}
F.~Calegari, S.~Garoufalidis and D.~Zagier, \emph{Bloch groups, algebraic $K$-theory, units, and Nahm's conjecture}, preprint (2017), arXiv:1712.04887.

\bibitem{CMP16}
C.~Calinescu, A.~Milas and M.~Penn, \emph{Vertex algebraic structure of principal subspaces of basic $A_{2n}^{(2)}$-modules}, J. Pure Appl. Algebra \textbf{220} (2016), 1752--1784.

\bibitem{CW25}
Z.~Cao and L.~Wang, \emph{Some new modular rank four Nahm sums as lift-dual of rank three examples}, preprint (2025), arXiv:2508.12468.

\bibitem{CW26}
Z.~Cao and L.~Wang, \emph{Some new modular Nahm sums of ranks 3 and 4}, preprint (2026), arXiv:2606.13590.

\bibitem{CTW26}
S.~Chern, C.~Tran and T.~Wakhare, \emph{Tadpole Nahm sum as a Wronskian}, preprint (2026), arXiv:2607.22904.

\bibitem{Cigler03}
J.~Cigler, \emph{A new class of $q$-Fibonacci polynomials}, Electron. J. Combin. \textbf{10} (2003), Research Paper 19, 15~pp. \href{https://doi.org/10.37236/1712}{doi:10.37236/1712}.

\bibitem{Singular}
W.~Decker, G.-M.~Greuel, G.~Pfister and H.~Sch\"onemann, \textsc{Singular} --- A computer algebra system for polynomial computations, \url{https://www.singular.uni-kl.de}.

\bibitem{Gollnitz67}
H.~G\"ollnitz, \emph{Partitionen mit Differenzenbedingungen}, J. Reine Angew. Math. \textbf{225} (1967), 154--190.

\bibitem{Gordon65}
B.~Gordon, \emph{Some continued fractions of the Rogers--Ramanujan type}, Duke Math. J. \textbf{32} (1965), 741--748.

\bibitem{Kanade19}
S.~Kanade, \emph{Searching for modular companions}, preprint (2019), arXiv:1902.05199.

\bibitem{Li26}
H.~Li, \emph{Proofs of five conjectural identities on modular rank four Nahm sums}, preprint (2026), arXiv:2606.25866.

\bibitem{MZ10}
J.~Mc~Laughlin and P.~Zimmer, \emph{General WP-Bailey chains}, Ramanujan J. \textbf{22} (2010), 11--31.

\bibitem{MW24}
A.~Milas and L.~Wang, \emph{Modularity of Nahm sums for the tadpole diagram}, Int. J. Number Theory \textbf{20} (2024), 73--101. \href{https://doi.org/10.1142/S1793042124500040}{doi:10.1142/S1793042124500040}; arXiv:2301.04532.

\bibitem{Nahm07}
W.~Nahm, \emph{Conformal field theory and torsion elements of the Bloch group}, in: Frontiers in Number Theory, Physics, and Geometry~II, Springer, Berlin, 2007, pp.~67--132.

\bibitem{SW26}
C.~Shi and L.~Wang, \emph{Modularity of tadpole Nahm sums in ranks 4 and 5}, J. Number Theory \textbf{284} (2026), 214--245. \href{https://doi.org/10.1016/j.jnt.2025.12.010}{doi:10.1016/j.jnt.2025.12.010}; arXiv:2504.17737.

\bibitem{SW26b}
C.~Shi and L.~Wang, \emph{Nahm sums dual to Zagier's rank three examples and related identities}, preprint (2026), arXiv:2607.23257v2.

\bibitem{VZ11}
M.~Vlasenko and S.~Zwegers, \emph{Nahm's conjecture: asymptotic computations and counterexamples}, Commun. Number Theory Phys. \textbf{5} (2011), 617--642. arXiv:1104.4008.

\bibitem{Zagier07}
D.~Zagier, \emph{The dilogarithm function}, in: Frontiers in Number Theory, Physics, and Geometry~II, Springer, Berlin, 2007, pp.~3--65.

\end{thebibliography}


\end{document}
